% ======================================================================
% CAP25 v13 insert --- body only, fifth-pass merged (chain-push +
% uploaded final + additional distinct-label full-petal chart) checked.  Designed for direct \input into
% cs25_cap_v12.tex immediately BEFORE \section{Discussion} (i.e. between
% \cref{sec:answers} and \cref{sec:discussion}).  All labels are
% prefixed v13- to avoid collisions with the v12 source.  Theorem
% environments, cleveref, enumitem, and booktabs are taken from the
% main file.  No bibliography items are added; all \cite keys below
% already exist in the v12 bibliography.
% ======================================================================
\providecommand{\F}{\mathbb F}
\providecommand{\RS}{\mathrm{RS}}
\providecommand{\ImgFib}{\operatorname{ImgFib}}
\providecommand{\Paid}{\operatorname{Paid}}
\providecommand{\Dloc}{\mathcal D}
\providecommand{\PP}{\mathbb P}
\providecommand{\E}{\mathbb E}
\providecommand{\Var}{\operatorname{Var}}
\providecommand{\Span}{\operatorname{span}}
\providecommand{\ceil}[1]{\left\lceil #1\right\rceil}
\providecommand{\floor}[1]{\left\lfloor #1\right\rfloor}

% ============================ BEGIN V13 BODY ============================

\section{v13 insertion map and verification status}\label{sec:v13-map}

This insert extends the certificate interface of \cref{sec:prize-program} and the ledger machinery of \cref{sec:discharge} by three groups of compilers: exact integer threshold compilers, list-side (L1) lower and upper compilers with a sunflower residual interface, and residue-line (M1) tools aimed at the aperiodic band of Open Problem~\ref{prob:band}.  Every statement below is either proved in full, proved by explicit citation of a v12 theorem, or is an exact-arithmetic certificate whose defining integer computation is printed.  Nothing in this insert claims to resolve \cref{prob:band}; its role is to make the future inputs of \cref{thm:conditional-mca,thm:conditional-list} auditable, to pay the strata that can already be paid, and to pin the residual branches by name.

\begin{center}
\small
\begin{tabular}{@{}p{0.24\linewidth}p{0.36\linewidth}p{0.33\linewidth}@{}}
\toprule
Block & What is proved here & How v13 uses it \\
\midrule
Threshold compilers (\cref{sec:v13-threshold}) & monotone integer staircases, endpoint conventions, lower/upper corridors, coarse steepness envelopes, budget windows, an \emph{exact} high-agreement tangent cell, status-valued quotient cells, a sharpened extension-pole conversion cell, and bounded quotient-census certificates & strengthens the certificate grammar of \cref{sec:prize-program}; replaces informal phrases such as ``adjacent crossing'' by exact integer lemmas; \cref{prop:onestep} becomes the one-step case of \cref{thm:v13-staircase,thm:v13-corridor} \\
\addlinespace
L1 list side (\cref{sec:v13-l1}) & planted quotient-core lower counts at arbitrary sub-period slack, complete polynomial-folding planted cores, exact list budget windows, a general $\kappa$-intersecting agreement bound matching \cref{thm:johnson-list}, and proved sunflower paid layers with distinct-label zero-excess and all-fixed-excess full-petal charts & adds lower-side list certificates and list residual interfaces; the primitive image-fiber theorem is \emph{not} claimed and is stated as \cref{prob:v13-primitive-image-fiber} \\
\addlinespace
M1 residue-line side (\cref{sec:v13-m1}) & exact first/second moments and an exact variance ledger for random split locators, common-GCD and quotient-pullback reductions, fixed-dimensional and projective-plane Conjecture-F bounds, quotient seam arithmetic, Hankel/exchange identities, and SPI deficiency-one eliminants with explicit degree caps & supplies a toolkit for the aperiodic branch of \cref{prob:band}; random-word and fixed-dimensional statements are \emph{not} promoted to a worst-case M1 local limit \\
\bottomrule
\end{tabular}
\end{center}

\begin{remark}[fourth-pass verification status]\label{rem:v13-verification}
Beyond the third-pass chain-push audit, every exact-arithmetic certificate in this insert has been independently recomputed, and the constructive witnesses have been checked by exhaustive small-field enumeration.  Specifically: the quotient-census tables of \cref{thm:v13-census} and the planted crossing tables of \cref{prop:v13-dyadic-planted} were recomputed by exact integer arithmetic, scanning \emph{all} admissible scales below each printed crossing, not only the printed predecessor; the rate-$1/2$ census identity is proved in closed form in \cref{lem:v13-census-identity}; the exact tangent witness of \cref{prop:v13-tangent} was verified by brute force over $\F_{13}$ \textup(the constructed pair has exactly $r+1$ support-wise MCA-bad finite slopes, matching \cref{thm:deep-mca}\textup); the planted quotient-core construction of \cref{thm:v13-planted} was verified over $\F_{17}$ at all slacks $1\le\sigma<M$; and the moment formulas of \cref{thm:v13-first-moment,thm:v13-second-moment} were verified by exhaustive enumeration of all word pairs over $\F_5$ and, locator-pair by locator-pair through the joint syndrome image, over $\F_7$, exercising every overlap rank class.  The high-agreement tangent cell is exact for Reed--Solomon rows \textup(\cref{prop:v13-tangent}\textup); the quotient cell is status-valued and cites the v12 support and image ledgers rather than restating them; the extension-pole compiler exposes its conversion certificate and uses the sharper collision denominator $m+\kappa(L-1)$ already implicit in the proof of \cref{cor:extension-pole-deep-list-floor}; the sunflower section proves its concrete layer injection, its residual-free distinct-label zero-excess full-petal chart strengthened to a union bound over missed cores, and a new residual-free distinct-label all-fixed-excess full-petal chart; repeated-label higher-excess full petals, mixed petals, and growing-defect classifications remain certificate or residual statements; the split-locator concentration statement uses the exact Johnson-neighborhood degree $D_t(n,j)$ and an exact variance ledger; and identically rank-deficient SPI pencils are separated as a named class \textup(0\textup) instead of being silently counted.
\end{remark}

\section{Threshold certificate compilers}\label{sec:v13-threshold}

\subsection{Integer staircases and corridors}

\begin{definition}[integer numerator staircase]\label{def:v13-staircase}
Fix an agreement interval
\[
        I=\{a_{\min},a_{\min}+1,\ldots,a_{\max}\}
\]
and a nonincreasing function
\[
        N:I\longrightarrow \mathbb Z_{\ge0}.
\]
For MCA or line-MCA rows, $N(a)$ is the number of bad sampled parameters at agreement at least $a$; the closed integer Hamming radius is $r=n-a$.  For list rows one may use the same agreement coordinate, so the list numerator is nonincreasing in $a$.  Given a denominator $Q$ and a target $\varepsilon^*$, put
\[
        B_*:=\floor{\varepsilon^* Q}.
\]
An agreement $a$ is \emph{safe} if $N(a)\le B_*$ and \emph{unsafe} if $N(a)>B_*$.
\end{definition}

\begin{theorem}[staircase localization]\label{thm:v13-staircase}
Let $N:I\to\mathbb Z_{\ge0}$ be nonincreasing.  Exactly one of the following holds:
\begin{enumerate}[label=\textup{(\roman*)}]
\item every $a\in I$ is safe;
\item every $a\in I$ is unsafe;
\item there is a unique first safe agreement $a_*\in I$ such that the safe set is $\{a\in I:a\ge a_*\}$ and the unsafe set is $\{a\in I:a<a_*\}$.
\end{enumerate}
In the interior case, when $a_*>a_{\min}$, one has
\[
        N(a_*-1)>B_*\ge N(a_*).
\]
\end{theorem}

\begin{proof}
The predicate $N(a)\le B_*$ is monotone nondecreasing in $a$: if it holds at $a$, then for $a'\ge a$ we have $N(a')\le N(a)\le B_*$.  A monotone Boolean function on a finite linearly ordered set is either identically false, identically true, or has a unique first true point.  The displayed inequalities are exactly the assertion that the point just before the first true point is false while the first true point is true.
\end{proof}

\begin{corollary}[closed-radius endpoint convention]\label{cor:v13-endpoint}
In the interior case of \cref{thm:v13-staircase}, the largest safe closed integer radius is
\[
        r_{\mathrm{safe}}=n-a_*,
\]
and the first unsafe closed integer radius above it is $r_{\mathrm{safe}}+1$.  As a real closed-ball threshold under the convention $a(\delta)=\ceil{(1-\delta)n}$ of \cref{sec:prize-program}, the supremum of safe radii is
\[
        \frac{n-a_*+1}{n},
\]
and this endpoint is not attained when $N(a_*-1)>B_*$.
\end{corollary}

\begin{proof}
The integer radius corresponding to agreement $a$ is $n-a$, so the first safe agreement $a_*$ gives the largest safe integer radius $n-a_*$.  A real radius $\delta$ is safe exactly when $\ceil{(1-\delta)n}\ge a_*$, which holds if and only if $(1-\delta)n>a_*-1$, i.e.\ for all $\delta<(n-a_*+1)/n$.  At $\delta=(n-a_*+1)/n$ the agreement requirement is $a_*-1$, which is unsafe in the interior case.
\end{proof}

\begin{theorem}[corridor from lower and upper certificates]\label{thm:v13-corridor}
Suppose that $L(a)\le N(a)\le U(a)$ are certified lower and upper bounds on the same staircase.  Define
\[
        A_{\mathrm{unsafe}}:=\{a\in I:L(a)>B_*\},
        \qquad
        A_{\mathrm{safe}}:=\{a\in I:U(a)\le B_*\}.
\]
If an interior first safe agreement $a_*$ exists, then
\[
        1+\max A_{\mathrm{unsafe}}\le a_*\le \min A_{\mathrm{safe}},
\]
with the conventions $1+\max\varnothing=a_{\min}$ and $\min\varnothing=a_{\max}$.
\end{theorem}

\begin{proof}
If $L(a)>B_*$, then $N(a)>B_*$, so $a$ is unsafe and the first safe point is strictly larger than $a$.  Taking the largest such certified unsafe point gives the lower bound.  If $U(a)\le B_*$, then $N(a)\le B_*$, so $a$ is safe and the first safe point is at most $a$.  Taking the smallest such certified safe point gives the upper bound.  The displayed conventions are exactly the empty-set cases.
\end{proof}

\begin{remark}[relation to the v12 interface]\label{rem:v13-staircase-v12}
\Cref{prop:onestep} is the one-step case of \cref{thm:v13-staircase,thm:v13-corridor}: a matched pair of certificates $L(a_0)>B_*\ge U(a_0+1)$ pins $a_*=a_0+1$ exactly.  In the conditional theorems \cref{thm:conditional-mca,thm:conditional-list}, the hypothesized band inputs enter precisely as the upper staircase $U$, and the unsafe floors of \cref{sec:pincer} enter as the lower staircase $L$; \cref{thm:v13-corridor} is the exact bookkeeping of the sandwich between them at a fixed budget.
\end{remark}

\begin{theorem}[coarse steepness compiler]\label{thm:v13-steepness}
Let $J\subseteq I$ be an interval of consecutive agreements, let $M(a)>0$ be a positive model count, and suppose a row packet proves
\[
        \frac{M(a)}{E}\le N(a)\le E M(a)\qquad(E\ge1)
\]
for every $a\in J$.  For a single adjacent pair $a,a+1\in J$, if
\[
        \frac{M(a)}{M(a+1)}>E^2,
\]
then the two ambiguity intervals
\[
        \left[\frac{M(a)}E,EM(a)\right]
        \quad\text{and}\quad
        \left[\frac{M(a+1)}E,EM(a+1)\right]
\]
are disjoint, and a fixed budget $B_*$ can lie in at most one of them.

If the same separation holds for every adjacent pair in $J$, then the ambiguity intervals are pairwise disjoint in their natural order.  Consequently at most one agreement in $J$ can remain undecided by the coarse envelope for a fixed budget $B_*$; every other agreement has a certified safe or unsafe sign.
\end{theorem}

\begin{proof}
For the adjacent pair, the displayed inequality is equivalent to
\[
        \frac{M(a)}E>EM(a+1),
\]
which is precisely the assertion that the lower end of the $a$-interval is above the upper end of the $(a+1)$-interval.  If $B_*<M(b)/E$, then $N(b)\ge M(b)/E>B_*$ and $b$ is certified unsafe.  If $B_*\ge EM(b)$, then $N(b)\le EM(b)\le B_*$ and $b$ is certified safe.  Thus only budgets inside the ambiguity interval for $b$ can be undecided.

When the adjacent separation holds throughout $J$, induction gives
\[
        \frac{M(a)}E>EM(a+1)\ge \frac{M(a+1)}E>EM(a+2)\ge\cdots
\]
for consecutive points, using $E\ge1$, so all ambiguity intervals are pairwise disjoint.  A single integer budget can belong to at most one member of a pairwise disjoint collection of intervals, and the sign criterion of the previous paragraph certifies all remaining agreements.
\end{proof}

\subsection{Paid ledger cells}

\begin{definition}[high-agreement tangent exact cell]\label{def:v13-tangent-cell}
For a Reed--Solomon row of length $n$ and dimension $k<n$, put $r(A)=n-A$ and
\[
        R_{\tan}:=\floor{\frac{n-k}{3}}.
\]
The high-agreement tangent paid cell is
\[
        \Paid_{\tan}^{\mathrm{hi}}(n,k,A):=r(A)+1
\]
when $0\le r(A)\le R_{\tan}$, and is unavailable outside this range unless a separate tangent/common-line theorem is supplied.
\end{definition}

\begin{proposition}[exact tangent cell]\label{prop:v13-tangent}
Let $C=\RS[F,D,k]$ with $|D|=n$ and $k<n$.  If $A=n-r$ and
\[
        3r\le n-k,
\]
then the maximum over received lines of the number of finite support-wise MCA-bad slopes at closed agreement threshold $A$ is exactly $r+1=n-A+1$, and the same maximum for no-loss CA-bad slopes is also exactly $r+1$.  Hence \cref{def:v13-tangent-cell} is an exact paid cell throughout the high-agreement range.
\end{proposition}

\begin{proof}
The upper bound is \cref{thm:deep-mca} applied at radius $r/n$: the Reed--Solomon minimum Hamming weight is $w=n-k+1$, so the hypothesis $3r\le w-1$ there is exactly $3r\le n-k$, and the theorem bounds the number of CA- or MCA-bad finite slopes of every received line by $r+1$.

For the matching lower witness, choose a set
\[
        T=\{t_0,\ldots,t_r\}\subseteq D
\]
of size $r+1$ \textup(possible since $r+1\le n$\textup) and choose distinct field elements $\gamma_0,\ldots,\gamma_r\in F$ \textup(possible since $|F|>n\ge r+1$ for $D\subseteq F^\times$, and $|F|\ge n\ge r+1$ in general\textup).  Define two received words by
\[
        f_2(t_i)=1,
        \qquad
        f_1(t_i)=-\gamma_i,
        \qquad
        f_1(x)=f_2(x)=0\quad (x\in D\setminus T).
\]
For the slope $\gamma_i$, the word $f_1+\gamma_i f_2$ vanishes at $t_i$ and off $T$, hence agrees with the zero codeword on
\[
        S_i:=D\setminus (T\setminus\{t_i\}),
        \qquad |S_i|=n-r=A .
\]

We first record the two degree inequalities used below.  If $r=0$, then $n-r-1=n-1\ge k$ and $n-2r-1=n-1\ge k$.  If $r\ge1$, then $r+1\le 3r\le n-k$ gives $n-r-1\ge k$, and $2r\le 3r-r\le n-k-r\le n-k-1$ gives $n-2r-1\ge k$.

\emph{Support-wise badness.}  Suppose $P_1,P_2\in F[X]_{<k}$ satisfy $P_1=f_1$ and $P_2=f_2$ on $S_i$.  On $D\setminus T\subseteq S_i$ both received words vanish, and $|D\setminus T|=n-r-1\ge k$, so each $P_j$ has at least $k$ roots and is the zero polynomial.  But $t_i\in S_i$ and $f_2(t_i)=1\ne0$, a contradiction.  Hence the pair $(f_1,f_2)$ is not simultaneously explained on $S_i$, while the point $f_1+\gamma_i f_2$ is explained there by the zero codeword: each $\gamma_i$ is support-wise MCA-bad at agreement threshold $A$, and the $r+1$ slopes $\gamma_i$ are distinct.

\emph{No-loss badness.}  Let $S\subseteq D$ be any support with $|S|\ge A=n-r$ and suppose $P_1,P_2\in F[X]_{<k}$ explain $(f_1,f_2)$ on $S$.  Since $|T|=r+1$, inclusion--exclusion gives $|S\cap T|\ge|S|+|T|-n\ge1$, and
\[
        |S\cap(D\setminus T)|\ge |S|-|T|\ge n-2r-1\ge k .
\]
Hence both $P_j$ vanish at $\ge k$ points and are identically zero, contradicting $f_2(t)=1$ at any point $t\in S\cap T$.  So no support of size $\ge A$ simultaneously explains the pair, i.e.\ the pair is no-loss far at radius $r/n$; since each $f_1+\gamma_if_2$ is explained on $S_i$, the same $r+1$ slopes are no-loss CA-bad.  Combining the witness with the upper bound of \cref{thm:deep-mca} proves exactness.
\end{proof}

\begin{remark}[tangent-cell calibration]\label{rem:v13-tangent-calibration}
The exactness makes the high-agreement staircase pointwise computable: in the tangent range the exact numerator is $N(A)=n-A+1$, which is strictly decreasing, so \cref{thm:v13-staircase} localizes with a one-step transition.  As a worked integer example, take the calibration row $n=512$, $k=256$ with sampling denominator $Q=17^{32}$ and target $2^{-128}$.  Exact arithmetic gives $B_*=\floor{17^{32}/2^{128}}=6$ \textup($17^{32}$ has $131$ bits\textup), and the tangent range is $r\le R_{\tan}=85$, i.e.\ $A\ge427$.  Inside this range, $N(A)=n-A+1\le6$ first holds at $a_*=507$: agreement $506$ is unsafe with $N=7>6$ and agreement $507$ is safe with $N=6\le6$.  The largest safe closed integer radius for the tangent-range staircase is $r_{\mathrm{safe}}=5$, and the real closed-ball supremum of \cref{cor:v13-endpoint} is $6/512$, not attained.  This is a calibration example for the compiler, not a deployed row.
\end{remark}

\begin{definition}[quotient paid status cell]\label{def:v13-quotient-status}
Let $\mathcal C$ be a finite set of declared quotient fiber sizes $c\mid n$ and let $A$ be an agreement threshold.  The conservative safe-sum support count is
\[
U_{\mathrm{sum}}(n,A,\mathcal C):=
\sum_{c\in\mathcal C}\sum_{B=A}^{n}
\binom{n/c}{\floor{B/c}}
\binom{n-c\floor{B/c}}{B-c\floor{B/c}},
\]
the quotient-remainder count of \cref{sec:discharge}: for each $c$ and each size $B$, choose $\floor{B/c}$ complete $c$-fibers among the $n/c$ fibers and then a residual set of size $B-c\floor{B/c}$ among the remaining $n-c\floor{B/c}$ points.  The quotient paid cell at $(n,A,\mathcal C)$ is status-valued:
\[
\Paid_{\mathrm{quot}}=
\begin{cases}
\mathrm{EXACT\_IMAGE}, & \text{an image-lcm certificate as in \cref{thm:exact-quotient-image-lcm-ledger}},\\
\mathrm{EXACT\_SUPPORT}, & \text{an exact block-profile union as in \cref{thm:exact-divisor-block-support-ledger}},\\
\mathrm{SAFE\_SUM}(U_{\mathrm{sum}}), & \text{otherwise}.
\end{cases}
\]
Unsafe quotient lower cells should carry zone tags, for example norm-exact, collision-open, or cap-floor.  A collision-open zone is interval-valued and must not be silently replaced by a point estimate.
\end{definition}

\begin{proposition}[safe-sum validity]\label{prop:v13-quotient-safe-sum}
Let $C=\RS[F,D,k]$ and $A\ge k$.  If a quotient packet declares that every finite bad slope in the cell is witnessed on at least one support counted by the quotient-remainder model above, then $U_{\mathrm{sum}}(n,A,\mathcal C)$ is a valid conservative upper numerator for the union of the declared quotient finite-slope cells at agreement at least $A$, for support-wise MCA-bad and for no-loss CA-bad slopes alike.
\end{proposition}

\begin{proof}
The displayed binomial product counts, for each declared $c$ and size $B$, the supports consisting of $\floor{B/c}$ complete $c$-fibers plus a residual set outside those fibers; summing over $B\ge A$ and $c\in\mathcal C$ is a union bound over the declared support family, which may double-count supports with several quotient descriptions but cannot undercount them.  By \cref{lem:one-support-one-line}, each fixed support pays at most one finite bad slope, in either badness sense.  The declared coverage hypothesis then converts the support union bound into a finite-slope numerator.  This is the safe-sum instantiation of \cref{prop:divisor-union-support-ledger}; when the exact union or the image-lcm coalescing is available, the sharper grades of \cref{def:v13-quotient-status} replace it.
\end{proof}

\begin{proposition}[extension-pole conversion cell]\label{prop:v13-extension}
Let $B\subseteq F$ be the generated field inside the line field, with $q_{\mathrm{gen}}=|B|$ and $q_{\mathrm{line}}=|F|$.  If $q_{\mathrm{gen}}=q_{\mathrm{line}}$, then there is no proper extension-only line-parameter cell:
\[
        \Paid_{\mathrm{ext}}^{\mathrm{only}}=0.
\]
Assume $q_{\mathrm{gen}}<q_{\mathrm{line}}$, put $m=q_{\mathrm{line}}-q_{\mathrm{gen}}$, and let $\mathcal P$ be a base list of size $L\ge1$.  Suppose a pole-conversion certificate assigns to every $P\in\mathcal P$ and every pole $\alpha\in F\setminus B$ a finite extension parameter $\lambda_\alpha(P)$ such that:
\begin{enumerate}[label=\textup{(\roman*)}]
\item for each fixed pole $\alpha$, distinct values among $\{\lambda_\alpha(P):P\in\mathcal P\}$ give distinct extension-only bad finite line parameters for the lifted row;
\item for every $P\ne P'$ in $\mathcal P$,
\[
        \#\{\alpha\in F\setminus B:\lambda_\alpha(P)=\lambda_\alpha(P')\}\le \kappa .
\]
\end{enumerate}
Then some pole produces at least
\[
        \mathrm{ExtPole}(q_{\mathrm{line}},q_{\mathrm{gen}},\kappa,L)
        :=
        \left\lceil
        \frac{L(q_{\mathrm{line}}-q_{\mathrm{gen}})}
        {q_{\mathrm{line}}-q_{\mathrm{gen}}+\kappa(L-1)}
        \right\rceil
\]
distinct extension-only bad finite line parameters.  When the conversion is evaluation at the pole, one may take $\kappa=k$ for a degree-$\le k$ auxiliary list and $\kappa=k-1$ for an ordinary degree-$<k$ Reed--Solomon list.  For safe-side extension charts over $B$, a closed parameter set of degree $\Delta$ and dimension $e$ contributes at most $\Delta q_{\mathrm{gen}}^{\,e}$ affine $B$-parameters, by \cref{thm:extension-line-dimension-degree-ledger}.
\end{proposition}

\begin{proof}
If $B=F$, then $F\setminus B=\varnothing$, so no line parameter is extension-only.  Now assume $m>0$.  For a pole $\alpha$, let $M_\alpha$ be the number of distinct values among $\{\lambda_\alpha(P):P\in\mathcal P\}$, let $n_{\alpha,\lambda}$ be the multiplicity of the value $\lambda$, and let $C_\alpha:=\sum_\lambda\binom{n_{\alpha,\lambda}}2$ be the number of colliding unordered pairs at $\alpha$.  Summing the pairwise certificate \textup{(ii)} over unordered pairs of list elements gives
\[
        \sum_{\alpha\in F\setminus B}C_\alpha
        \le \kappa\binom L2 ,
\]
so some pole $\alpha$ has $C_\alpha\le \kappa\binom L2/m$.  At this pole, Cauchy--Schwarz gives
\[
        L^2
        =\Bigl(\sum_\lambda n_{\alpha,\lambda}\Bigr)^2
        \le
        M_\alpha\sum_\lambda n_{\alpha,\lambda}^2
        =M_\alpha(L+2C_\alpha)
        \le
        M_\alpha\Bigl(L+\frac{\kappa L(L-1)}m\Bigr),
\]
hence
\[
        M_\alpha\ge \frac{Lm}{m+\kappa(L-1)} .
\]
Since $M_\alpha$ is an integer, the displayed ceiling is certified, and condition \textup{(i)} converts the distinct pole values into distinct extension-only bad line parameters.  For the evaluation conversion, two distinct polynomials of degree $\le k$ \textup(resp.\ $<k$\textup) agree at most $k$ \textup(resp.\ $k-1$\textup) times on $F\setminus B$, giving the stated $\kappa$.  The safe-side clause is \cref{thm:extension-line-dimension-degree-ledger} verbatim.
\end{proof}

\begin{remark}[relation to the v12 extension-pole floors]\label{rem:v13-extension-v12}
\Cref{prop:punctured-simple-pole-floor} and \cref{cor:extension-pole-deep-list-floor} realize the conversion certificate concretely: the conversion is evaluation of the base list at a simple pole $\alpha\in F\setminus B$, condition \textup{(i)} holds with genuinely extension-valued witness lines, and condition \textup{(ii)} holds with $\kappa=k$ there.  The proof of \cref{cor:extension-pole-deep-list-floor} internally derives the denominator $m+k(L-1)$ before weakening it to $m+kL$ for display; \cref{prop:v13-extension} promotes the sharper denominator to the official compiler value.  The older denominator $q_{\mathrm{line}}-q_{\mathrm{gen}}+\kappa L$ remains a valid conservative fallback and may be printed when only the weaker collision certificate is available.
\end{remark}

\subsection{Budget windows and bounded quotient census}

\begin{theorem}[budget windows and dodge selection]\label{thm:v13-windows}
Let $N$ be the actual integer numerator for a fixed object, and suppose a certificate gives $L\le N\le K$.  Put $B(q)=\floor{q/2^{128}}$.  Then
\[
        B(q)<L \Rightarrow \text{certified unsafe},
        \qquad
        B(q)\ge K \Rightarrow \text{certified safe for this object},
\]
and the only unresolved budgets are
\[
        L\le B(q)<K.
\]
Equivalently, the unresolved field-size interval is
\[
        L2^{128}\le q\le K2^{128}-1.
\]
If the lower proof is $L=K-g$, then the largest tolerable missing mass while still proving unsafety is
\[
        g_{\max}=K-B(q)-1.
\]
\end{theorem}

\begin{proof}
The safety condition for the fixed object is $N\le B(q)$.  If $B(q)<L\le N$, then $N>B(q)$ and the object certifies unsafety.  If $N\le K\le B(q)$, then the object is safe under this certificate.  The complement is precisely $L\le B(q)<K$.  The floor definition turns this double inequality into the displayed field-size interval.  Finally, if $L=K-g$, unsafety requires $K-g>B(q)$, equivalently $g<K-B(q)$, and the largest such integer is $K-B(q)-1$.
\end{proof}

For a $2$-power quotient order $N'$, let $n_1=N'/2$.  The characteristic-zero antipodal quotient count used by the quotient census of \cite{Cho26b} is
\[
        A_2(N',\ell')
        :=
        \sum_{\substack{u\ge0,\ t=\ell'-2u\ge0\\ u\le n_1-t}}
        \binom{n_1}{t}2^t .
\]

\begin{lemma}[rate-$1/2$ census identity]\label{lem:v13-census-identity}
For every even $N'\ge2$, with $n_1=N'/2$,
\[
        A_2(N',n_1+1)=\frac{3^{n_1}-1}{2}.
\]
\end{lemma}

\begin{proof}
At $\ell'=n_1+1$ the summation constraints reduce to a single parity class.  Indeed $t=n_1+1-2u$, so $t\equiv n_1+1\pmod2$; the constraint $u\le n_1-t$ becomes $u\le n_1-(n_1+1-2u)=2u-1$, i.e.\ $u\ge1$, which is equivalent to $t\le n_1-1$; and $t\ge0$ closes the range.  Hence
\[
        A_2(N',n_1+1)
        =\sum_{\substack{0\le t\le n_1-1\\ t\equiv n_1+1\ (\mathrm{mod}\ 2)}}\binom{n_1}t2^t
        =\sum_{\substack{0\le t\le n_1\\ t\not\equiv n_1\ (\mathrm{mod}\ 2)}}\binom{n_1}t2^t,
\]
where the last step notes that $t=n_1$ has the excluded parity, and $t=n_1+1$ is out of range.  From $\sum_t\binom{n_1}t2^t=3^{n_1}$ and $\sum_t\binom{n_1}t(-2)^t=(-1)^{n_1}$, the odd-$t$ class sums to $(3^{n_1}-(-1)^{n_1})/2$ and the even-$t$ class to $(3^{n_1}+(-1)^{n_1})/2$.  If $n_1$ is even the class $t\not\equiv n_1$ is the odd class and $(-1)^{n_1}=1$; if $n_1$ is odd it is the even class and $(-1)^{n_1}=-1$.  In both cases the sum is $(3^{n_1}-1)/2$.
\end{proof}

\begin{theorem}[bounded quotient census, exact arithmetic certificate]\label{thm:v13-census}
Let $B_{\max}=2^{128}-1$.  In the relaxed even-scale model, where $N'$ ranges over the even integers with $\rho N'\in\mathbb Z$, the first admissible scales at which $A_2(N',\rho N'+1)>B_{\max}$ are
\[
\begin{array}{c|cccc}
\rho & 1/2 & 1/4 & 1/8 & 1/16\\ \hline
N'_{\mathrm{relaxed}} & 164 & 176 & 248 & 384 .
\end{array}
\]
For dyadic quotient orders $N'=2^a$, the corresponding first scales are
\[
\begin{array}{c|cccc}
\rho & 1/2 & 1/4 & 1/8 & 1/16\\ \hline
N'_{\mathrm{dyadic}} & 256 & 256 & 256 & 512 .
\end{array}
\]
Thus the exact integer quotient endgame at the $2^{-128}$ budget sees only bounded quotient scales, independent of the deployed blocklength, once the row has been reduced to this canonical quotient-count model.
\end{theorem}

\begin{proof}
The count $A_2(N',\ell')$ is an explicit finite integer sum, so the statement is an exact arithmetic certificate.  For each fixed $\rho$, the certificate evaluates $A_2(N',\rho N'+1)$ by exact integer arithmetic at \emph{every} admissible scale up to and including the printed first crossing, confirms $A_2\le B_{\max}$ at every admissible scale strictly below the crossing, and confirms $A_2>B_{\max}$ at the crossing; no floating-point estimate and no unproved monotonicity is used.  Writing $\operatorname{bits}(X)=\floor{\log_2X}+1$, so that $\operatorname{bits}(X)\le128\iff X\le B_{\max}$, the immediate predecessor and crossing values are
\[
\begin{array}{c|cc|cc}
\rho & N'_{\mathrm{prev}} & \operatorname{bits} A_2 & N'_{\mathrm{first}} & \operatorname{bits} A_2\\ \hline
1/2  & 162 & 128 & 164 & 129\\
1/4  & 172 & 127 & 176 & 130\\
1/8  & 240 & 127 & 248 & 131\\
1/16 & 368 & 124 & 384 & 129
\end{array}
\qquad
\begin{array}{c|cc|cc}
\rho & N'_{\mathrm{prev}} & \operatorname{bits} A_2 & N'_{\mathrm{first}} & \operatorname{bits} A_2\\ \hline
1/2  & 128 & 101 & 256 & 202\\
1/4  & 128 &  95 & 256 & 190\\
1/8  & 128 &  68 & 256 & 135\\
1/16 & 256 &  87 & 512 & 172
\end{array}
\]
for the relaxed even model \textup(left\textup) and the dyadic model \textup(right\textup).  At rate $1/2$ the scanned values agree with the closed form of \cref{lem:v13-census-identity}, which is strictly increasing in $N'$, so the rate-$1/2$ rows also follow from the identity alone.
\end{proof}

\begin{definition}[dyadic planted quotient profile]\label{def:v13-qprofile}
For $n=2^m$, $k=\rho n$, and integer slack $\sigma=\ceil{\eta n}$, define
\[
K_Q(n,k,\sigma):=
\max_{\substack{M=2^e\mid\gcd(n,k)\\ \sigma<M\\ k/M\le n/M-1}}
\binom{n/M-1}{k/M},
\]
with value $0$ if the active set is empty.
\end{definition}

\begin{proposition}[bounded active quotient order]\label{prop:v13-qprofile}
If $\sigma\ge n/256$, then every active scale in $K_Q$ has quotient order $N=n/M<256$; because $N$ is dyadic in \cref{def:v13-qprofile}, this means $N\le128$.  If $\sigma\ge n/128$, then every active scale has $N<128$, hence dyadically $N\le64$.
\end{proposition}

\begin{proof}
Activity requires $\sigma<M=n/N$, hence $N<n/\sigma$.  The two displayed strict bounds follow from the two lower bounds on $\sigma$, and the dyadic constraint improves the strict inequalities to $N\le128$ and $N\le64$ respectively.
\end{proof}

\begin{proposition}[polynomial residual absorption]\label{prop:v13-integrality}
Fix a candidate scale and suppose a future theorem bounds an unpaid residual contribution by
\[
        n^C\Phi(n,q)
\]
for some explicit proxy $\Phi$.  If the verified entropy calculation gives
\[
        -\log_2\bigl(n^C\Phi(n,q)\bigr)>0,
\]
then the residual integer contribution is zero.
\end{proposition}

\begin{proof}
The residual contribution is a nonnegative integer bounded above by a real number strictly smaller than $1$.  The only such integer is $0$.
\end{proof}

\section{L1 list-side compilers and sunflower residuals}\label{sec:v13-l1}

Let $D\subset F$ be an evaluation domain and let $C=\RS[F,D,k]$ be a Reed--Solomon code of length $n=|D|$.  For a received word $U:D\to F$ and an agreement threshold $s$, write
\[
        \ImgFib_U^{<k}(s):=
        \{P\in F[X]_{<k}: \#\{x\in D:U(x)=P(x)\}\ge s\},
\]
and, for an auxiliary degree bound $d$ and a subdomain $T\subseteq D$,
\[
        \ImgFib_W^{\le d}(s;T):=
        \{G\in F[X]_{\le d}: \#\{x\in T:W(x)=G(x)\}\ge s\}.
\]
These are the actual list objects used below; they count codewords rather than raw supports.

\subsection{Planted quotient-core lower counts}

\begin{theorem}[planted quotient-core lower count]\label{thm:v13-planted}
Let $D=\alpha H\subset F^\times$ be a multiplicative coset of a cyclic subgroup $H$ of order $n$, let $C=\RS[F,D,k]$, and suppose $M\mid\gcd(n,k)$ satisfies
\[
        1\le \sigma<M,
        \qquad
        k/M\le n/M-1.
\]
Then there is a received word $U:D\to F$ with at least
\[
        P_M(n,k):=\binom{n/M-1}{k/M}
\]
distinct codewords agreeing with $U$ on at least $k+\sigma$ positions.  Consequently
\[
\max_U \bigl|\ImgFib_U^{<k}(k+\sigma)\bigr|
\ge
\max_{\substack{M\mid\gcd(n,k)\\ \sigma<M\\ k/M\le n/M-1}}
\binom{n/M-1}{k/M}.
\]
\end{theorem}

\begin{proof}
Let $K\le H$ be the unique subgroup of order $M$ and let $Q:=\{x^M:x\in D\}$ be the quotient image, so that the map $\pi:D\to Q$, $\pi(x)=x^M$, has fibers equal to the $K$-cosets inside $D$, each of size exactly $M$, and $|Q|=N:=n/M$.  Choose one quotient value $b_0\in Q$ and a set $T\subseteq\pi^{-1}(b_0)$ of size $\sigma$; this is possible since $\sigma<M$.  Put
\[
        E_T(X):=\prod_{t\in T}(X-t).
\]
For every subset $A\subseteq Q\setminus\{b_0\}$ of size $k/M$, define
\[
        L_A(X):=E_T(X)\prod_{b\in A}(X^M-b).
\]
Every point of the fiber $\pi^{-1}(b)$ with $b\in A$ is a root of $X^M-b$, and every point of $T$ is a root of $E_T$; these point sets are disjoint since $b_0\notin A$.  Hence $L_A$ vanishes at at least $M\cdot(k/M)+\sigma=k+\sigma$ points of $D$.

Since
\[
        \prod_{b\in A}(X^M-b)=X^k+\text{terms of degree at most }k-M,
\]
and $\deg E_T=\sigma<M$, the part of $L_A$ of degree at least $k$ is independent of $A$, namely
\[
        U_0(X):=X^kE_T(X).
\]
Thus $L_A=U_0+R_A$ with $\deg R_A<k$.  Let $U$ be the received word $x\mapsto U_0(x)$ on $D$ and let $P_A:=-R_A\in F[X]_{<k}$.  For every root $x\in D$ of $L_A$,
\[
        0=L_A(x)=U(x)-P_A(x),
\]
so $P_A$ agrees with $U$ on at least $k+\sigma$ positions.

It remains to check distinctness.  If $P_A$ and $P_{A'}$ give the same codeword on $D$, then $R_A-R_{A'}$ has degree $<k<n$ and vanishes on all $n$ points of $D$, so $R_A=R_{A'}$ as polynomials, hence $L_A=L_{A'}$, and dividing by the fixed nonzero polynomial $E_T$ gives
\[
        \prod_{b\in A}(X^M-b)=\prod_{b\in A'}(X^M-b).
\]
The substitution homomorphism $F[Y]\to F[X]$, $G(Y)\mapsto G(X^M)$, is injective \textup(the monomials $X^{Mi}$ are linearly independent\textup), so $\prod_{b\in A}(Y-b)=\prod_{b\in A'}(Y-b)$ in $F[Y]$, and unique factorization into the distinct monic linear factors gives $A=A'$.  \textup(This argument needs no separability of $X^M-b$ and is valid in every characteristic.\textup)  There are $\binom{N-1}{k/M}$ choices of $A$, which proves the claim.
\end{proof}

\begin{corollary}[complete polynomial-folding planted core]\label{cor:v13-polyfold-planted}
Let $H\subset F$ be finite, let $\phi\in F[X]$ have degree $M>0$, and assume the map $\phi:H\to Q$ is onto with every fiber $\phi^{-1}(b)\cap H$ of size exactly $M$.  Suppose
\[
        M\mid k,
        \qquad
        1\le\sigma<M,
        \qquad
        k/M\le |Q|-1 .
\]
Then there is a received word $U:H\to F$ with at least
\[
        \binom{|Q|-1}{k/M}
\]
distinct degree-$<k$ Reed--Solomon codewords agreeing with $U$ on at least $k+\sigma$ positions.
\end{corollary}

\begin{proof}
Choose $b_0\in Q$ and a subset $T\subseteq \phi^{-1}(b_0)\cap H$ of size $\sigma$.  Put
\[
        E_T(X):=\prod_{t\in T}(X-t),
        \qquad
        \ell:=k/M.
\]
For every $A\subseteq Q\setminus\{b_0\}$ of size $\ell$, define
\[
        L_A(X):=E_T(X)\prod_{b\in A}(\phi(X)-b).
\]
Its roots in $H$ include the $\sigma$ planted points of $T$ and the $M\ell=k$ points of the complete fibers over $A$, disjointly, hence at least $k+\sigma$ points.  Since
\[
        \prod_{b\in A}(Y-b)=Y^\ell+\text{terms of degree at most }\ell-1,
\]
we have
\[
        L_A(X)=E_T(X)\phi(X)^\ell+R_A(X),
        \qquad
        \deg R_A\le \sigma+M(\ell-1)<k.
\]
Let $U$ be the word $x\mapsto E_T(x)\phi(x)^\ell$ on $H$ and put $P_A:=-R_A\in F[X]_{<k}$; then $U=P_A$ at every root of $L_A$ in $H$.

For distinctness, note $|H|=M|Q|$ by the complete-fiber hypothesis, and $k=M\ell\le M(|Q|-1)<|H|$.  If $P_A$ and $P_{A'}$ define the same codeword on $H$, then $R_A-R_{A'}$ has degree $<k<|H|$ and vanishes on $H$, hence $R_A=R_{A'}$ and
\[
        \prod_{b\in A}(\phi(X)-b)=\prod_{b\in A'}(\phi(X)-b).
\]
The substitution map $F[Y]\to F[X]$, $G(Y)\mapsto G(\phi(X))$, is injective because $\deg\phi>0$, so $\prod_{b\in A}(Y-b)=\prod_{b\in A'}(Y-b)$ and $A=A'$.  The number of choices is $\binom{|Q|-1}{\ell}$.
\end{proof}

\begin{remark}[relation to the v12 fiber lemmas and extension persistence]\label{rem:v13-planted-v12}
\Cref{thm:v13-planted} is the exact-count planted form of the quotient-core construction of \cite{Cho26b}: a single received word carries the whole binomial sublist, with no pigeonhole loss and no subfield denominator.  By contrast, \cref{lem:fiber} pigeonholes locator subsets over their slope values to extract one heavy fiber at slack two; the planted form is the right tool on the list side, where no slope average is available.  \Cref{cor:v13-polyfold-planted} is the map-smooth analogue: its complete-uniform-fiber hypothesis is exactly the condition of \cref{def:map-smooth} used by \cref{lem:phi-fiber}, so the planted core applies verbatim to Chebyshev and circle-derived foldings.  Finally, all these list floors persist over every field extension: $\RS[B,D,k]\subseteq\RS[F,D,k]$ for $D\subseteq B\subseteq F$, so the planted sublist survives unchanged when the row is lifted, exactly as in the list-side persistence discussion of \cref{sec:subfield}.
\end{remark}

\begin{corollary}[list unsafe test]\label{cor:v13-list-unsafe}
For a list denominator $q_{\mathrm{list}}$ and target $2^{-128}$, the planted quotient-core packet proves a list row unsafe whenever
\[
\max_M\binom{n/M-1}{k/M}>
\floor{\frac{q_{\mathrm{list}}}{2^{128}}},
\]
where the maximum ranges over the active divisors of \cref{thm:v13-planted}.
\end{corollary}

\begin{proof}
By \cref{thm:v13-planted}, the maximum list size at agreement $k+\sigma$ is at least the displayed planted count.  If that count exceeds the allowed numerator $\floor{q_{\mathrm{list}}/2^{128}}$, the normalized list quantity exceeds $2^{-128}$.
\end{proof}

\begin{proposition}[dyadic planted crossings]\label{prop:v13-dyadic-planted}
For $n=2^m$, $k=\rho n$, and
\[
        \rho\in\{1/2,\,1/4,\,1/8,\,1/16\},
\]
the first dyadic quotient orders $N=n/M$ at which the planted count exceeds $2^{128}-1$ are
\[
\begin{array}{c|cccc}
\rho & 1/2 & 1/4 & 1/8 & 1/16\\ \hline
N & 256 & 256 & 256 & 512 ,
\end{array}
\]
and the exact bit lengths of $\binom{N-1}{\rho N}$ at the predecessor and first crossing are
\[
\begin{array}{c|cc|cc}
\rho & N_{\mathrm{prev}} & \operatorname{bits} & N_{\mathrm{first}} & \operatorname{bits}\\ \hline
1/2  & 128 & 124 & 256 & 251\\
1/4  & 128 & 100 & 256 & 204\\
1/8  & 128 &  67 & 256 & 136\\
1/16 & 256 &  83 & 512 & 169.
\end{array}
\]
\end{proposition}

\begin{proof}
The planted count at quotient order $N$ is $\binom{N-1}{\rho N}$, and $\operatorname{bits}(X)\le128\iff X\le2^{128}-1$; the displayed bit lengths were checked by exact integer arithmetic on binomial coefficients, not floating-point logarithms.  It remains to justify that the predecessor check identifies the first dyadic crossing.  Along dyadic orders the counts are monotone: from a subset $S\subseteq\{1,\ldots,N-1\}$ of size $\rho N$, form a subset of $\{1,\ldots,2N-1\}$ of size $2\rho N$ by adjoining a fixed set of $\rho N$ elements of $\{N,\ldots,2N-1\}$.  This map is injective, so
\[
        \binom{N-1}{\rho N}\le \binom{2N-1}{2\rho N}.
\]
Hence once the predecessor is below the budget, every smaller dyadic order is below it as well.
\end{proof}

\begin{corollary}[planted list budget windows]\label{cor:v13-list-windows}
If a planted exact count is $P$ and the certified lower count is $L\le P$, then the only unresolved list budgets are
\[
        L\le \floor{q_{\mathrm{list}}/2^{128}}<P,
\]
equivalently $L2^{128}\le q_{\mathrm{list}}\le P2^{128}-1$.
\end{corollary}

\begin{proof}
This is \cref{thm:v13-windows} with $q=q_{\mathrm{list}}$, $K=P$.
\end{proof}

\subsection{Agreement-threshold Johnson bounds}

\begin{theorem}[$\kappa$-intersecting agreement bound]\label{thm:v13-johnson-list}
Let $\Omega$ be a finite set of size $n_0$, let $W:\Omega\to F$ be a received word, and let $\mathcal F$ be a family of words $\Omega\to F$ such that any two distinct members of $\mathcal F$ agree on at most $\kappa$ points of $\Omega$.  For an agreement threshold $s$ with $0<s\le n_0$, let $L$ be the number of members of $\mathcal F$ agreeing with $W$ on at least $s$ points.
\begin{enumerate}[label=\textup{(\roman*)}]
\item If $2s>n_0+\kappa$, then $L\le1$.
\item If $s^2>\kappa n_0$, then
\[
        L
        \le
        \frac{n_0(s-\kappa)}{s^2-\kappa n_0}.
\]
\end{enumerate}
In particular, for $C=\RS[F,\Omega,k]$ and $\kappa=k-1$,
\[
        \bigl|\ImgFib_W^{<k}(s)\bigr|\le\frac{n_0(s-k+1)}{s^2-(k-1)n_0}
        \qquad\text{when }s^2>(k-1)n_0,
\]
and for the auxiliary degree-$\le d$ code on a petal domain, $\kappa=d$.
\end{theorem}

\begin{proof}
Let $P_1,\ldots,P_L$ be the listed members and $A_i:=\{x\in\Omega:P_i(x)=W(x)\}$, so $|A_i|\ge s$ and, for $i\ne j$, $|A_i\cap A_j|\le\kappa$ by the intersecting hypothesis \textup(two words agreeing with $W$ at a point agree with each other there\textup).

\textup{(i)}  For $L\ge2$, inclusion--exclusion gives $|A_1\cap A_2|\ge2s-n_0>\kappa$, a contradiction.

\textup{(ii)}  Assume $L\ge1$ and put $m_x:=\#\{i:x\in A_i\}$ and $I:=\sum_{x\in\Omega}m_x\ge Ls$.  Counting incident pairs,
\[
        \sum_{x\in\Omega}\binom{m_x}2=\sum_{i<j}|A_i\cap A_j|\le\binom L2\kappa .
\]
First suppose $Ls\le n_0$.  Then $L\le n_0/s$, and since $s\le n_0$,
\[
        \frac{n_0}s\le\frac{n_0(s-\kappa)}{s^2-\kappa n_0}
        \iff
        s^2-\kappa n_0\le s^2-\kappa s
        \iff
        \kappa s\le\kappa n_0,
\]
which holds; note $s-\kappa>0$ because $s^2>\kappa n_0\ge\kappa s$.  Now suppose $Ls>n_0$.  The function $\varphi(y)=y(y-1)$ is convex and nondecreasing on $[1,\infty)$, so Jensen and $I/n_0\ge Ls/n_0>1$ give
\[
        n_0\cdot\frac{Ls}{n_0}\Bigl(\frac{Ls}{n_0}-1\Bigr)
        \le
        n_0\,\varphi\Bigl(\frac I{n_0}\Bigr)
        \le
        \sum_x m_x(m_x-1)
        \le
        L(L-1)\kappa .
\]
Multiplying out, $Ls(Ls-n_0)\le n_0L(L-1)\kappa$, hence $s(Ls-n_0)\le n_0(L-1)\kappa$, i.e.
\[
        L(s^2-\kappa n_0)\le n_0(s-\kappa),
\]
which is the claim.
\end{proof}

\begin{remark}[relation to the v12 Johnson theorem]\label{rem:v13-johnson-v12}
For $\kappa=k-1$ this is exactly the $m=1$ case of \cref{thm:johnson-list}, restated in agreement-threshold list coordinates; the numerator $n_0(s-k+1)$ is the sharp one and strictly improves the numerator $n_0(n_0-k+1)$ that a plain Cauchy--Schwarz argument yields.  The abstract $\kappa$-intersecting form is recorded because the sunflower reduction below applies it with $\kappa=d$ to an auxiliary code that is not a row of the challenge.
\end{remark}

\subsection{Sunflower and petal interfaces}

The statements in this subsection are compiler reductions: they become theorems about an actual sunflower decomposition only when the decomposition supplies the stated layer maps.  This is the safe v13 form: the paid pieces are proved, while the mixed-petal and growing-defect classification remains a named residual branch.  The first proposition records one concrete situation in which the required layer injection is automatic.

\begin{definition}[concrete sunflower core layer]\label{def:v13-concrete-sunflower}
Let $\Omega\subset F$ be a finite evaluation set for $\RS[F,\Omega,k]$.  Fix pairwise disjoint subsets
\[
        Y,\quad R_0,\quad T_1,\ldots,T_M\subseteq\Omega,
        \qquad |Y|=k-1,
        \qquad |T_i|=\ell=\sigma+1,
\]
a reference polynomial $P_\star\in F[X]_{<k}$, and labels $c_1,\ldots,c_M\in F$.  Put
\[
        L_Y(X)=\prod_{y\in Y}(X-y),
\]
and assume the received word $U:\Omega\to F$ satisfies
\[
        U=P_\star\quad\text{on }Y\cup R_0,
        \qquad
        U=P_\star+c_iL_Y\quad\text{on }T_i .
\]
For $D_0\subseteq Y$ write $d=|D_0|$, $L_{D_0}(X)=\prod_{y\in D_0}(X-y)$, and $T:=\bigcup_iT_i$.  The concrete $(D_0,R_0)$ layer consists of those $P\in F[X]_{<k}$ that agree with $U$ on at least $k+\sigma$ points of $\Omega$, agree with $U$ on all of $(Y\setminus D_0)\cup R_0$, and have no agreements with $U$ outside $(Y\setminus D_0)\cup R_0\cup T$.
\end{definition}

\begin{proposition}[concrete sunflower layer injection]\label{prop:v13-concrete-sunflower}
In the setting of \cref{def:v13-concrete-sunflower}, the map
\[
        P\longmapsto
        G_P:=\frac{P-P_\star}{L_{Y\setminus D_0}}
\]
is well-defined and injects the concrete $(D_0,R_0)$ layer into $F[X]_{\le d}$.  Moreover $G_P$ agrees on at least
\[
        a=\sigma+d+1-|R_0|
\]
points of the petal domain $T$ with the auxiliary word
\[
        U_{D_0}(x):=c_iL_{D_0}(x),\qquad x\in T_i,
\]
and if the layer retains the residual set $R_0$, then in addition $G_P(x)=0$ for every $x\in R_0$.
\end{proposition}

\begin{proof}
Since $P$ and $P_\star$ agree on $Y\setminus D_0$ \textup(both equal $U=P_\star$ there\textup), the difference $P-P_\star$ is divisible by $L_{Y\setminus D_0}$, which has degree $k-1-d$; the quotient $G_P$ therefore has degree at most $(k-1)-(k-1-d)=d$.  The quotient determines $P=P_\star+G_PL_{Y\setminus D_0}$, so the map is injective.

The layer has at least $k+\sigma$ agreements in total, of which at most $|Y\setminus D_0|+|R_0|=(k-1-d)+|R_0|$ lie off the petal domain, by the containment condition of the definition.  Hence at least
\[
        (k+\sigma)-(k-1-d)-|R_0|=\sigma+d+1-|R_0|=a
\]
agreement points lie in $T$.  If $x\in T_i$ is such a point, then $x\notin Y$, so $L_{Y\setminus D_0}(x)\ne0$, and
\[
        G_P(x)L_{Y\setminus D_0}(x)
        =P(x)-P_\star(x)
        =U(x)-P_\star(x)
        =c_iL_Y(x)
        =c_iL_{D_0}(x)L_{Y\setminus D_0}(x);
\]
cancelling the nonzero factor gives $G_P(x)=c_iL_{D_0}(x)=U_{D_0}(x)$.  Finally, if $x\in R_0$, then $P(x)=U(x)=P_\star(x)$ and $x\notin Y$, so $G_P(x)=0$.
\end{proof}

\begin{definition}[fixed sunflower auxiliary layer]\label{def:v13-sunflower-layer}
Fix a missed-core set $D_0$ of size $d$ with monic locator $L_{D_0}(X)=\prod_{y\in D_0}(X-y)$, a retained residual subset $R_0$ of size $r$, pairwise disjoint petals $T_1,\ldots,T_M$ of common size $\sigma+1$ disjoint from $D_0\cup R_0$, and labels $c_i\in F$.  Let $T:=\bigcup_iT_i$ and define
\[
        U_{D_0}(x):=c_iL_{D_0}(x),\qquad x\in T_i.
\]
A fixed $(D_0,R_0)$ sunflower layer is a finite set $\mathcal L(D_0,R_0)$ together with an injective map
\[
        \Psi:\mathcal L(D_0,R_0)\hookrightarrow F[X]_{\le d}
\]
such that every $G=\Psi(P)$ agrees with $U_{D_0}$ on at least
\[
        a:=\sigma+d+1-r
\]
points of $T$.  Extra conditions forcing zeros on $R_0$ are called retained zero constraints.
\end{definition}

\begin{proposition}[sunflower core-defect reduction]\label{prop:v13-pma}
For every fixed sunflower auxiliary layer in the sense of \cref{def:v13-sunflower-layer},
\[
        |\mathcal L(D_0,R_0)|
        \le
        \bigl|\ImgFib_{U_{D_0}}^{\le d}(a;T)\bigr|,
\]
the auxiliary degree-$\le d$ list on the petal domain $T$ at agreement $a=\sigma+d+1-r$.  Dropping the retained zero constraints on $R_0$ only enlarges the right-hand side, so it is safe for upper bounds.
\end{proposition}

\begin{proof}
By definition, $\Psi$ injects the fixed layer into the set of degree-$\le d$ polynomials agreeing with $U_{D_0}$ on at least $a$ points of $T$, which is the displayed auxiliary list.  Adding zero constraints cuts out a subset of it; removing them can only increase the counted set.
\end{proof}

\begin{corollary}[few-petal Johnson boundary]\label{cor:v13-pma-johnson}
In the notation of \cref{prop:v13-pma}, the auxiliary petal domain has size $|T|=M(\sigma+1)$, and two distinct degree-$\le d$ polynomials agree on at most $d$ points, so \cref{thm:v13-johnson-list} applies with $\kappa=d$.  If $a^2>d|T|$, the fixed $(D_0,R_0)$ auxiliary list has size at most
\[
        \frac{|T|\,(a-d)}{a^2-d\,|T|}.
\]
Equivalently, for $d>0$, the few-petal Johnson-covered range is exactly
\[
        M\le
        \floor{\frac{a^2-1}{d(\sigma+1)}}.
\]
\end{corollary}

\begin{proof}
The bound is \cref{thm:v13-johnson-list}\textup{(ii)} with $n_0=|T|$, $s=a$, $\kappa=d$ \textup(if $a>|T|$ the list is empty and the bound is vacuous\textup).  Since all quantities are integers, the strict inequality $a^2>M(\sigma+1)d$ is equivalent to $M(\sigma+1)d\le a^2-1$, which is the stated floor.
\end{proof}

\begin{definition}[full-petal fixed-excess counting chart]\label{def:v13-full-petal-chart}
A full-petal counting chart with cofactor excess $e=d-\ell$ is a parameterization of a full-petal layer with the following certificate; the certificate is for the whole layer being charged, and if several independent charts are used the bounds below must be summed over the chart index.
\begin{enumerate}[label=\textup{(\roman*)}]
\item the zero-excess layer $e=0$ is determined by an unordered pair of petals and one scalar in $F$;
\item for every $e\ge1$, the layer is determined by a subset of the $M$ petals and at most $e+1$ scalars in $F$.
\end{enumerate}
\end{definition}

\begin{proposition}[distinct-label zero-excess full-petal chart]\label{prop:v13-zero-excess-chart}
Work in a concrete sunflower layer \textup(\cref{def:v13-concrete-sunflower}\textup) with $R_0=\varnothing$, petal size $\ell=\sigma+1$, and $d=\ell$, and assume the labels $c_1,\ldots,c_M$ are pairwise distinct.  Consider a residual-free full-petal sublayer: every listed codeword $P$ agrees with $U$ on the full union of a set $I(P)$ of complete petals, and the listedness inequality forces $|I(P)|\ge2$.  Then the map
\[
        P\longmapsto(\{i,j\},\alpha_i),
\]
where $i<j$ are the two smallest indices in $I(P)$ and $\alpha_i$ is the scalar constructed in the proof, is injective, so the sublayer has size at most
\[
        \binom M2\,|F| .
\]
Moreover the injection recovers the missed-core locator $L_{D_0}$ itself, hence $D_0$; consequently the same bound $\binom M2|F|$ holds for the \emph{union} of these sublayers over all missed-core sets $D_0\subseteq Y$ of size $\ell$.  In particular the chart supplies the zero-excess certificate of \cref{def:v13-full-petal-chart}.
\end{proposition}

\begin{proof}
Write $L:=L_{D_0}$ and let $L_i$ be the monic locator of the petal $T_i$; both are monic of degree $\ell$.  For a codeword $P$ in the sublayer write, as in \cref{prop:v13-concrete-sunflower},
\[
        P-P_\star=G\,L_{Y\setminus D_0},
        \qquad \deg G\le d=\ell .
\]
If $i\in I(P)$, then $P$ agrees with $U=P_\star+c_iL_Y$ at every point of $T_i$, so by \cref{prop:v13-concrete-sunflower} the polynomial $G-c_iL$, of degree at most $\ell$, vanishes at all $\ell$ points of $T_i$; hence
\[
        G-c_iL=\alpha_iL_i
        \qquad\text{for some }\alpha_i\in F .
\]
Subtracting the identities for the two smallest touched petals $i<j$,
\[
        (c_j-c_i)L=\alpha_iL_i-\alpha_jL_j .
\]
Comparing leading coefficients of the three monic degree-$\ell$ locators gives
\[
        c_j-c_i=\alpha_i-\alpha_j,
        \qquad\text{so}\qquad
        \alpha_j=\alpha_i-(c_j-c_i).
\]
Given the pair $\{i,j\}$ and the scalar $\alpha_i$, the label difference $c_j-c_i\ne0$ determines $\alpha_j$, then
\[
        L=\frac{\alpha_iL_i-\alpha_jL_j}{c_j-c_i}
\]
determines $L$, then $G=c_iL+\alpha_iL_i$, and finally $P=P_\star+G\,L_{Y\setminus D_0}$, where $D_0$ is recovered as the root set of the monic squarefree polynomial $L$ inside $Y$ and $L_{Y\setminus D_0}=L_Y/L$.  Thus the map is injective, even jointly over the choice of $D_0$, and there are at most $\binom M2|F|$ values.
\end{proof}

\begin{proposition}[distinct-label fixed-excess full-petal chart]\label{prop:v13-distinct-excess-chart}
Work in a concrete sunflower layer \textup(\cref{def:v13-concrete-sunflower}\textup) with $R_0=\varnothing$, petal size $\ell=\sigma+1$, and pairwise distinct labels $c_1,\ldots,c_M$.  Fix $e\ge1$ and put $d=\ell+e$.  Consider the charged pairs $(P,D_0)$, where $D_0\subseteq Y$ has size $d$ and $P$ lies in a residual-free all-full sublayer with missed core $D_0$: there is a touched set $I(P,D_0)\subseteq[M]$ such that $P$ agrees with $U$ on every point of $T_i$ for $i\in I(P,D_0)$, has no petal agreements for $i\notin I(P,D_0)$, and has no agreements outside $(Y\setminus D_0)\cup\bigcup_{i\in I(P,D_0)}T_i$.  Then, for this fixed $e$, the number of charged pairs is at most
\[
        2^M |F|^{e+1}.
\]
Consequently the union of these residual-free all-full distinct-label sublayers over all missed cores $D_0$ also has size at most $2^M|F|^{e+1}$.  Equivalently, in this distinct-label residual-free all-full regime the higher-excess certificate of \cref{def:v13-full-petal-chart} is automatic, not an extra hypothesis.
\end{proposition}

\begin{proof}
Let $(P,D_0)$ be such a charged pair, put $L:=L_{D_0}$, and let $L_i$ be the monic locator of the petal $T_i$.  As in \cref{prop:v13-concrete-sunflower}, write
\[
        P-P_\star=G\,L_{Y\setminus D_0},
        \qquad \deg G\le d=\ell+e .
\]
If $i\in I:=I(P,D_0)$, full-petal agreement on $T_i$ gives
\[
        G-c_iL=L_iB_i
\]
for a uniquely determined polynomial $B_i$ of degree at most $e$, because the left side has degree at most $\ell+e$ and vanishes on the $\ell$ points of $T_i$.

The listedness threshold is $k+\sigma=k+\ell-1$.  Since the sublayer is residual-free and all-full, all agreements are contained in the $(k-1-d)$ retained core points and the $|I|\ell$ full petal points.  Hence
\[
        (k-1-d)+|I|\ell\ge k+\ell-1,
        \qquad\text{so}\qquad
        (|I|-1)\ell\ge d .
\]
In particular $I$ has at least two elements.  Let $i_0$ be the least element of $I$.  For every $j\in I\setminus\{i_0\}$, subtracting the full-petal identities gives
\[
        (c_{i_0}-c_j)L+L_{i_0}B_{i_0}=L_jB_j,
\]
so, since $c_{i_0}\ne c_j$,
\[
        L\equiv -(c_{i_0}-c_j)^{-1}L_{i_0}B_{i_0}\pmod {L_j}.
\]
The petal locators are pairwise coprime, so the touched set $I$ and the single polynomial $B_{i_0}$ determine, by the Chinese remainder theorem, the residue class of $L$ modulo
\[
        Q_I:=\prod_{j\in I\setminus\{i_0\}}L_j,
        \qquad \deg Q_I=(|I|-1)\ell\ge d .
\]
There is at most one monic polynomial of degree $d$ in this residue class: if $\deg Q_I>d$ this is immediate because the residue representative has degree $<\deg Q_I$, and if $\deg Q_I=d$ the difference of two monic degree-$d$ representatives has degree $<d$ while being divisible by $Q_I$.  Thus $L=L_{D_0}$ is recovered from $(I,B_{i_0})$.

Once $L$ is known, the actual charged pair recovers $D_0$ as the root set of the monic squarefree locator $L$ inside $Y$, then recovers $L_{Y\setminus D_0}=L_Y/L$, then recovers
\[
        G=c_{i_0}L+L_{i_0}B_{i_0},
        \qquad
        P=P_\star+G L_{Y\setminus D_0}.
\]
Hence the map $(P,D_0)\mapsto (I,B_{i_0})$ is injective.  There are at most $2^M$ possible touched sets and at most $|F|^{e+1}$ possible polynomials $B_{i_0}$ of degree at most $e$, proving the bound.
\end{proof}

\begin{remark}[higher-excess full petals after the distinct-label chart]\label{rem:v13-higher-excess-cert}
\Cref{prop:v13-distinct-excess-chart} pays residual-free all-full higher-excess petals when the petal labels are pairwise distinct.  For repeated labels, retained residual sets, partial or mixed petals, or charts that do not satisfy the residual-free all-full containment hypothesis, Chinese-remainder reconstructions may still certify the $2^Mq^{e+1}$ chart of \cref{def:v13-full-petal-chart}, but the reconstruction data must be printed as a certificate.  The insert does not infer those remaining higher-excess charts from listedness alone.
\end{remark}

\begin{proposition}[fixed-excess full-petal compiler]\label{prop:v13-fixed-excess}
For every fixed $E\ge0$, any single certified full-petal atlas satisfying \cref{def:v13-full-petal-chart} has total size at most
\[
        \binom M2 q + 2^M\sum_{e=1}^{E}q^{e+1},
        \qquad q=|F| ,
\]
the certified layer bounds being $\binom M2q$ at $e=0$ and $2^Mq^{e+1}$ at each $e\ge1$.
\end{proposition}

\begin{proof}
For $e=0$ the certificate chooses an unordered pair of petals and one scalar: at most $\binom M2q$ possibilities.  For fixed $e\ge1$ it chooses a subset of the $M$ petals in at most $2^M$ ways and $e+1$ scalars in at most $q^{e+1}$ ways.  In the residual-free pairwise-distinct-label all-full case, \cref{prop:v13-distinct-excess-chart} proves exactly this certificate; in other cases it remains a printed atlas certificate.  Summing over $1\le e\le E$ proves the claim.
\end{proof}

\begin{corollary}[fixed-excess polynomiality]\label{cor:v13-fixed-excess-poly}
At the intended L1 lower cutoff, suppose $M=O(\log n)$, $q\le n^Q$ for a fixed $Q$, and $E$ is fixed.  Then the fixed-excess full-petal contribution certified by \cref{prop:v13-fixed-excess} is polynomial in $n$.
\end{corollary}

\begin{proof}
The term $\binom M2q$ is polynomial because $M=O(\log n)$ and $q\le n^Q$.  For fixed $e\le E$, the term $2^Mq^{e+1}$ is $n^{O(1)}\cdot n^{Q(e+1)}$, hence polynomial, and a fixed finite sum of polynomially bounded terms is polynomial.
\end{proof}

\begin{remark}[L1 evidence is not a theorem]\label{rem:v13-l1-evidence}
The small-field auxiliary scans and worst-word sweeps that motivated the sunflower split, together with the explicit growing-defect full-petal examples, should be retained only as evidence and route cuts.  They support the residual decomposition but do not prove polynomial growth for the mixed-petal or growing-defect branches.
\end{remark}

\begin{problem}[primitive image-fiber theorem]\label{prob:v13-primitive-image-fiber}
Prove a polynomial bound for the stabilizer-primitive part of $\ImgFib_U^{<k}(k+\sigma)$ once quotient-periodic codewords are charged separately and the generated-field entropy reserve and lower cutoff clear.
\end{problem}

\begin{problem}[mixed-petal and growing-defect residuals]\label{prob:v13-l1-residuals}
Classify or bound mixed-petal sunflower extras in the sub-Johnson region after the few-petal Johnson layer \textup(\cref{cor:v13-pma-johnson}\textup), the distinct-label all-full fixed-excess layer \textup(\cref{prop:v13-distinct-excess-chart}\textup), and any certified fixed-excess full-petal layers \textup(\cref{prop:v13-fixed-excess}\textup) are paid.  Determine whether the full-petal branch with $d-\ell\to\infty$ remains polynomial or produces a new obstruction ledger.
\end{problem}

\section{M1 residue-line tools and Conjecture-F reductions}\label{sec:v13-m1}

The safe-side program asks for upper bounds on bad line parameters after tangent, quotient, and extension cells have been charged.  In exact-agreement notation,
\[
        j=n-A,
        \qquad
        t=A-k,
        \qquad
        j+t=n-k,
\]
and the M1 problem is to control aperiodic residue-line packings in this $(j,t)$ window --- the residual kernel of \cref{prob:band}.  The results below are reusable local algebra and combinatorics; they do not prove the full worst-case M1 local limit.

\subsection{Split-locator moment calculus}

Let $F=\F_q$, let $D\subset F^\times$ have size $n$, and let $k<A\le n$.  Put $t=A-k\ge1$ and $j=n-A$.  For $R\subseteq D$ with $|R|=j$, define
\[
        \ell_R(X)=\prod_{r\in R}(X-r)
\]
and the locator-syndrome vector of a word $w:D\to F$,
\[
        S_R(w)=
        \Bigl(\sum_{x\in D}w(x)\,\ell_R(x)\,x^m\Bigr)_{m=0}^{t-1}
        \in F^t .
\]
The shifted convention $m=1,\ldots,t$ is equivalent on $D\subset F^\times$: it is the composite of $S_R$ with the linear automorphism $w(x)\mapsto xw(x)$ of $F^D$, which preserves the uniform measure.  For independent uniform words $u,v:D\to F$, call $R$ \emph{aligned} if $S_R(v)\ne0$ and $S_R(u)\in F\,S_R(v)$, and let $N_A(u,v)$ be the number of aligned locators.

\begin{theorem}[exact first moment]\label{thm:v13-first-moment}
\[
        \E N_A
        =
        \binom{n}{j}\,(1-q^{-t})\,q^{1-t}.
\]
\end{theorem}

\begin{proof}
Fix $R$ and put $E:=D\setminus R$, $|E|=n-j=A$.  The linear map $w\mapsto S_R(w)$ has matrix entries $\ell_R(x)x^m$; the columns indexed by $x\in R$ vanish, and on $E$ we have $\ell_R(x)\ne0$, so scaling the column of $x$ by $\ell_R(x)^{-1}$ leaves the $t\times A$ Vandermonde block $(x^m)_{0\le m<t,\ x\in E}$ on distinct points.  Since $A=k+t\ge t$, any $t$ of its columns form an invertible Vandermonde matrix, so the map is surjective onto $F^t$.  A surjective linear image of a uniform vector is uniform, and $u,v$ are independent, so $(S_R(u),S_R(v))$ is uniform on $F^t\times F^t$.  The number of pairs $(a,b)$ with $b\ne0$ and $a\in Fb$ is $(q^t-1)q$, giving alignment probability $(q^t-1)q/q^{2t}=(1-q^{-t})q^{1-t}$; summing over the $\binom nj$ locators proves the formula.
\end{proof}

\begin{theorem}[two-locator joint rank]\label{thm:v13-joint-rank}
For $R,T\subseteq D$ of size $j$, put $c=|R\cap T|$.  The linear map
\[
        w\longmapsto (S_R(w),S_T(w))\in F^t\times F^t
\]
has rank
\[
        2t-\max(0,\,t-j+c).
\]
\end{theorem}

\begin{proof}
A linear relation among the $2t$ coordinate functionals is a pair of polynomials $A_0,B_0\in F[X]_{<t}$ such that $\sum_x w(x)\bigl(\ell_R(x)A_0(x)+\ell_T(x)B_0(x)\bigr)=0$ for all $w$, i.e.\ such that
\[
        \ell_R(X)A_0(X)+\ell_T(X)B_0(X)
\]
vanishes on all of $D$.  Its degree is at most $j+t-1\le n-1<n$, so it is the zero polynomial.  Let $G=\gcd(\ell_R,\ell_T)$, the locator of $R\cap T$, and write $\ell_R=Gr$, $\ell_T=Gs$ with $\deg G=c$, $\deg r=\deg s=j-c$, and $\gcd(r,s)=1$.  The identity becomes $rA_0+sB_0=0$, so $s\mid A_0$ and $r\mid B_0$: $A_0=sH$, $B_0=-rH$ for a polynomial $H$ with
\[
        \deg H\le t-1-(j-c).
\]
Hence the relation space has dimension $\max(0,\,t-j+c)$, and the rank is $2t$ minus this nullity.
\end{proof}

\begin{corollary}[covariance support]\label{cor:v13-covariance}
The alignment indicators for $R$ and $T$ are independent whenever
\[
        d(R,T):=j-|R\cap T|\ge t .
\]
All covariance is supported in the radius-$(t-1)$ Johnson neighborhood.
\end{corollary}

\begin{proof}
The condition $d(R,T)\ge t$ is $t-j+|R\cap T|\le0$, so by \cref{thm:v13-joint-rank} the map $w\mapsto(S_R(w),S_T(w))$ is surjective onto $F^{2t}$; hence $S_R(w)$ and $S_T(w)$ are independent uniform vectors for each of the independent words $u,v$, and the four syndrome vectors are mutually independent.  The alignment event for $R$ depends only on $(S_R(u),S_R(v))$ and that for $T$ only on $(S_T(u),S_T(v))$, so the two events are independent.
\end{proof}

\begin{theorem}[exact second moment]\label{thm:v13-second-moment}
Put $h(c):=\max(0,\,t-j+c)$ and, for $0\le h\le t$,
\[
P_h:=
\frac{q\,(q^h-1)\,q^{2(t-h)}+q^2\,(q^{t-h}-1)^2}{q^{4t-2h}} .
\]
Then
\[
\E N_A^2
=
\sum_{c=\max(0,2j-n)}^{j}
\binom{n}{j}\binom{j}{c}\binom{n-j}{j-c}\,
P_{h(c)} .
\]
\end{theorem}

\begin{proof}
Fix an ordered pair of locators $(R,T)$ with $|R\cap T|=c$ and put $h=h(c)$; the number of such ordered pairs is $\binom nj\binom jc\binom{n-j}{j-c}$, and $c$ ranges over $\max(0,2j-n)\le c\le j$.  It suffices to show that the probability that both $R$ and $T$ are aligned equals $P_h$.

\emph{Normal form of the joint image.}  Let $V\subseteq F^t\oplus F^t$ be the image of $w\mapsto(S_R(w),S_T(w))$; by \cref{thm:v13-joint-rank}, $\dim V=2t-h$, and both coordinate projections $\pi_1,\pi_2:V\to F^t$ are surjective \textup(each single-locator map is surjective, as in the first-moment proof\textup).  Put
\[
        W_1:=\{a\in F^t:(a,0)\in V\},
        \qquad
        W_2:=\{b\in F^t:(0,b)\in V\}.
\]
Since $\pi_1$ is surjective with kernel $\{0\}\oplus W_2$, we get $\dim W_2=(2t-h)-t=t-h$, and symmetrically $\dim W_1=t-h$.  The correspondence $a+W_1\mapsto b+W_2$ for $(a,b)\in V$ is a well-defined linear isomorphism $F^t/W_1\to F^t/W_2$ between $h$-dimensional spaces: well-defined because $(a,b),(a,b')\in V$ force $b-b'\in W_2$, injective and surjective by the symmetric statement and the dimension count.  Choose splittings $F^t=Z_i\oplus W_i$ with $\dim Z_i=h$, identify $Z_1\cong Z_2\cong F^h$ through the isomorphism, and correct the $W$-components by a linear section of $V$ over $Z_1$.  In the resulting coordinates,
\[
        V=U_h:=\{(z,p;\,z,q):z\in F^h,\ p,q\in F^{t-h}\},
\]
with the first block representing $S_R$ and the second $S_T$; alignment is invariant under these blockwise invertible changes of coordinates because each acts by an invertible linear map on the corresponding syndrome vector, preserving the relation $S(u)\in F\,S(v)$ and the condition $S(v)\ne0$.

\emph{Counting.}  The pair of words $(u,v)$ maps to a uniform element of $U_h\times U_h$, of size $q^{2(2t-h)}=q^{4t-2h}$.  Write the syndrome of $v$ as $(z,p;z,p')$.  If $z\ne0$ \textup($q^h-1$ choices, with $p,p'$ free in $q^{2(t-h)}$ ways\textup), then both $S_R(v)$ and $S_T(v)$ are nonzero, and simultaneous alignment of $u$ forces a common scalar: $S_R(u)=\lambda_R S_R(v)$ and $S_T(u)=\lambda_T S_T(v)$ share the first block $z_u=\lambda_Rz=\lambda_Tz$, so $\lambda_R=\lambda_T=:\lambda$, and each $\lambda\in F$ gives exactly one admissible syndrome of $u$, which lies in $U_h$; this contributes $q(q^h-1)q^{2(t-h)}$ aligned pairs.  If $z=0$, alignment requires $p\ne0$ and $p'\ne0$ \textup($(q^{t-h}-1)^2$ choices\textup), the two scalars are unconstrained and each choice $(\lambda_R,\lambda_T)\in F^2$ gives a distinct admissible syndrome of $u$ inside $U_h$; this contributes $q^2(q^{t-h}-1)^2$.  Dividing the total by $q^{4t-2h}$ gives $P_h$.
\end{proof}

\begin{corollary}[exact variance ledger]\label{cor:v13-exact-variance}
Let $p:=(1-q^{-t})q^{1-t}$.  Then
\[
\Var(N_A)=
\sum_{c=\max(0,2j-n)}^{j}
\binom{n}{j}\binom{j}{c}\binom{n-j}{j-c}
\bigl(P_{h(c)}-p^2\bigr),
\]
and the terms with $j-c\ge t$ vanish.
\end{corollary}

\begin{proof}
By \cref{thm:v13-first-moment} each locator is aligned with probability $p$, and the Vandermonde identity $\sum_c\binom jc\binom{n-j}{j-c}=\binom nj$ decomposes $(\E N_A)^2=\binom nj^2p^2$ over the same ordered overlap classes as \cref{thm:v13-second-moment}; subtracting termwise gives the display.  If $j-c\ge t$, then $h(c)=0$ and $P_0=q^2(q^t-1)^2/q^{4t}=p^2$, so those covariance terms are zero, in accordance with \cref{cor:v13-covariance}.
\end{proof}

\begin{corollary}[dependency-degree concentration]\label{cor:v13-dependency}
Let
\[
        D_t(n,j):=
        \sum_{d=0}^{\min(t-1,\,j,\,n-j)}
        \binom jd\binom{n-j}{d}
\]
be the size of the closed radius-$(t-1)$ Johnson ball.  Then
\[
        \Var(N_A)\le \E N_A\cdot D_t(n,j),
\]
so $\E N_A\ge D_t(n,j)\,n^s$ implies $\Var(N_A)/(\E N_A)^2\le n^{-s}$.  The coarser sufficient condition $\E N_A\ge t\,n^{2(t-1)+s}$ implies the same conclusion.
\end{corollary}

\begin{proof}
Write $N_A=\sum_RI_R$ with alignment indicators $I_R$.  By \cref{cor:v13-covariance}, $\operatorname{Cov}(I_R,I_T)=0$ unless $d(R,T)\le t-1$; the number of such $T$ for fixed $R$ is $\sum_{d\le t-1}\binom jd\binom{n-j}d=D_t(n,j)$ \textup(choose the $d$ points of $R$ to remove and the $d$ points outside to add\textup).  For every pair, $\operatorname{Cov}(I_R,I_T)\le\E(I_RI_T)\le\E I_R$.  Hence
\[
        \Var(N_A)=\sum_{R,T}\operatorname{Cov}(I_R,I_T)
        \le \sum_R D_t(n,j)\,\E I_R
        =\E N_A\cdot D_t(n,j),
\]
and dividing by $(\E N_A)^2$ gives the exact condition.  The coarse condition follows from $D_t(n,j)\le\sum_{d=0}^{t-1}n^{2d}\le t\,n^{2(t-1)}$.
\end{proof}

\begin{remark}[random-word scope]\label{rem:v13-moment-scope}
The moment calculus is an average over random word pairs.  It does not bound worst-case pairs; those require the aperiodic M1 local limit of \cref{prob:band}, a Conjecture-F incidence input, or exchange-rigidity input.
\end{remark}

\subsection{Conjecture-F reductions}

Let $K$ be a field, let $H\subset K$ be finite, and put $P_H(X)=\prod_{h\in H}(X-h)$.  Let $\Dloc_j(H)$ be the set of monic squarefree degree-$j$ divisors of $P_H$, i.e.\ of locators $L_S(X)=\prod_{h\in S}(X-h)$ with $S\subseteq H$, $|S|=j$.  Conjecture-F asks for upper bounds on $|\PP(W)\cap\Dloc_j(H)|$ for linear subspaces $W\le K[X]_{\le j}$; the lemmas below pay its reducible strata.

\begin{lemma}[common-GCD reduction]\label{lem:v13-gcd}
Let $W\le K[X]_{\le j}$ be a nonzero vector space, let $P=\PP(W)$, and let $E=P\cap\Dloc_j(H)$.  If every locator in $E$ is divisible by a fixed monic divisor $G$ of $P_H$ with $\deg G=w$, then division by $G$ injects $E$ into
\[
        \Dloc_{j-w}\bigl(H\setminus Z(G)\bigr)
\]
inside a projective linear space of dimension at most $\dim P$.
\end{lemma}

\begin{proof}
Division by the fixed polynomial $G$ is a linear injection from the subspace $W\cap GK[X]$ into $K[X]_{\le j-w}$; write $W_G$ for its image, so $\dim\PP(W_G)\le\dim P$.  If $[L]\in E$ with $L$ its monic representative, then $L/G$ is monic of degree $j-w$ and, since $L$ is squarefree with roots in $H$ and contains all roots of $G$, the quotient $L/G$ is squarefree with roots exactly $Z(L)\setminus Z(G)\subseteq H\setminus Z(G)$.  Thus $[L]\mapsto[L/G]$ is the claimed injection into $\PP(W_G)\cap\Dloc_{j-w}(H\setminus Z(G))$.
\end{proof}

\begin{lemma}[quotient-pullback recursion]\label{lem:v13-quot-pullback}
Assume $H=\mu_n$ and $\operatorname{char}K\nmid n$.  If $M\mid\gcd(n,j)$, then
\[
        g(Y)\longmapsto g(X^M)
\]
maps $\Dloc_{j/M}(\mu_{n/M})$ bijectively onto the $M$-periodic stratum of $\Dloc_j(\mu_n)$, i.e.\ onto the locators whose root sets are unions of complete $\mu_M$-cosets.  Intersecting with a projective linear space descends to an intersection with a projective linear space of no larger dimension in the quotient locator space.
\end{lemma}

\begin{proof}
Because $\operatorname{char}K\nmid n$, the polynomial $X^n-1$ is squarefree, and the $M$-power map $\mu_n\to\mu_{n/M}$ is surjective with all fibers of size $M$.  A monic squarefree divisor $g$ of $Y^{n/M}-1$ of degree $j/M$ pulls back to $g(X^M)$, a monic divisor of $X^n-1$ of degree $j$ whose root set is the union of the fibers over $Z(g)$, hence $M$-periodic and squarefree.  Conversely, an $M$-periodic squarefree divisor $L$ of $X^n-1$ has root set a union of fibers; its image $S\subseteq\mu_{n/M}$ has size $j/M$, and $g:=\prod_{b\in S}(Y-b)$ satisfies $g(X^M)=L$ \textup(both are monic with the same roots inside the squarefree $X^n-1$\textup).  This proves the bijection.  For the descent, the space $\{g:g(X^M)\in W\}$ is linearly isomorphic to $W\cap K[X^M]$ by injectivity of the substitution, hence has dimension at most $\dim W$.
\end{proof}

\begin{lemma}[dimension-one voting]\label{lem:v13-dim1}
Assume $j\ge1$.  Let $P=\PP(\Span(L_0,L_1))$ be a projective line in $K[X]_{\le j}$, and assume no $h\in H$ is a common zero of $L_0$ and $L_1$.  Then
\[
        |P\cap\Dloc_j(H)|\le\floor{\frac{|H|}{j}}.
\]
\end{lemma}

\begin{proof}
For each $h\in H$, the functional $[a:b]\mapsto aL_0(h)+bL_1(h)$ is nonzero by the common-zero hypothesis, so it has a unique kernel point on the projective line; say that $h$ votes for this point.  A locator point $[aL_0+bL_1]\in P\cap\Dloc_j(H)$ has exactly $j$ roots in $H$ and receives a vote precisely from each of them.  There are $|H|$ votes in total, so the number of locator points is at most $|H|/j$, hence at most the floor.
\end{proof}

\begin{lemma}[hyperplane concurrency]\label{lem:v13-concurrency}
Let $W\le K[X]_{\le j}$ be nonzero and assume no $h\in H$ is a common root of all elements of $W$.  For each $h$, let
\[
        E_h:=\{[L]\in\PP(W):L(h)=0\},
\]
a projective hyperplane of $\PP(W)$.  Then $\PP(W)\cap\Dloc_j(H)$ is exactly the set of projective points lying on at least $j$ of the hyperplanes $E_h$.
\end{lemma}

\begin{proof}
If $[L]\in\PP(W)\cap\Dloc_j(H)$, then $L$ has exactly $j$ distinct roots in $H$, so $[L]$ lies on the corresponding $j$ hyperplanes.  Conversely, if a nonzero $L\in W$ vanishes at $j$ or more distinct points of $H$, then $\deg L\le j$ forces $L$ to have exactly $j$ roots, all simple and all among those points; after monic normalization, $L\in\Dloc_j(H)$.
\end{proof}

\begin{theorem}[projective-plane pair bound]\label{thm:v13-dim2}
Under the hypotheses of \cref{lem:v13-concurrency}, assume $j\ge2$ and $\dim\PP(W)=2$.  Then
\[
        |\PP(W)\cap\Dloc_j(H)|
        \le
        \frac{\binom{|H|}{2}}{j-1},
\]
and if the evaluation lines $E_h$ are pairwise distinct, the sharper bound with denominator $\binom j2$ holds.
\end{theorem}

\begin{proof}
Here the hyperplanes $E_h$ are lines in the projective plane $\PP(W)$; distinct lines meet in exactly one point.  Group the $E_h$ into coincidence classes.  No class has multiplicity $\ge j$: if $E_{h_1}=\cdots=E_{h_j}$ for distinct $h_1,\ldots,h_j$, this common line is a two-dimensional projective space of polynomials of degree $\le j$ vanishing at the $j$ fixed points $h_1,\ldots,h_j$; but that vanishing space is spanned by the single locator $\prod_i(X-h_i)$, hence at most zero-dimensional projectively --- a contradiction.

Now let $[L]$ be a locator point and let $m_1,\ldots,m_s$ be the multiplicities of the line classes through $[L]$ among the $j$ lines $E_h$, $h\in Z(L)$, so $\sum_i m_i=j$ and each $m_i\le j-1$; in particular $s\ge2$.  Call a pair $\{h,h'\}\subseteq Z(L)$ a \emph{cross pair} if $E_h\ne E_{h'}$.  Fixing any one class, of multiplicity $m$, the number of cross pairs at $[L]$ is at least $m(j-m)$, which for integers $1\le m\le j-1$ is minimized at the endpoints with value $j-1$.  Each cross pair $\{h,h'\}$ determines the unique intersection point $E_h\cap E_{h'}$, so it can be charged to at most one locator point.  Since there are at most $\binom{|H|}2$ pairs of points of $H$ in total, the number of locator points is at most $\binom{|H|}2/(j-1)$.  When all $E_h$ are pairwise distinct, every locator point is incident to at least $j$ distinct lines, hence carries at least $\binom j2$ cross pairs, giving the sharper denominator.
\end{proof}

\begin{theorem}[fixed-dimensional Conjecture-F bound]\label{thm:v13-fixeddim}
Let $W\le K[X]_{\le j}$ have vector-space dimension $d+1$ with $0\le d\le j$, and assume $W$ is gcd-trivial on $H$, i.e.\ no point of $H$ is a common root of all elements of $W$.  Then
\[
        |\PP(W)\cap\Dloc_j(H)|\le \binom{|H|}{d}.
\]
More generally, if $C_0\subseteq H$ is the common root set of $W$ on $H$ and $d\le j-|C_0|$, then
\[
        |\PP(W)\cap\Dloc_j(H)|\le \binom{|H|-|C_0|}{d}.
\]
\end{theorem}

\begin{proof}
First assume gcd-triviality and fix a total order on $H$.  For a locator point $[L]$ with root set $R\subseteq H$, $|R|=j$, the vanishing subspace
\[
        W(-R):=\{G\in W:G|_R=0\}
\]
equals $K\cdot L$: it contains $L$, and any $G\in W\subseteq K[X]_{\le j}$ vanishing at the $j$ points of $R$ is a scalar multiple of the degree-$j$ locator of $R$.

Process the roots of $R$ in the fixed order, imposing one vanishing condition at a time, and record the roots at which the dimension strictly drops; call this set $S_0\subseteq R$.  At every stage the current subspace equals $W(-R')$ for the set $R'$ of processed roots, and a skipped root leaves the subspace unchanged, so by induction the current subspace equals $W(-S')$ for the set $S'$ of \emph{kept} roots so far; in particular $W(-S_0)=W(-R)=K\cdot L$ after all of $R$ is processed.  Each vanishing condition drops the dimension by at most one, and the total drop is $(d+1)-1=d$, so $|S_0|=d$ exactly.  The assignment $[L]\mapsto S_0$ is therefore a map into the $d$-subsets of $H$, and it is injective because $S_0$ determines $W(-S_0)=K\cdot L$, i.e.\ the point $[L]$.  This gives the bound $\binom{|H|}d$.

For the general version, every element of $W$ vanishes on $C_0$ and has degree $\le j$, so $W\subseteq G_{C_0}K[X]$ for $G_{C_0}:=\prod_{h\in C_0}(X-h)$; moreover every locator in the intersection vanishes on $C_0$, hence is divisible by $G_{C_0}$.  By \cref{lem:v13-gcd}, division by $G_{C_0}$ injects the intersection into a gcd-trivial instance on $H\setminus C_0$ with locator degree $j-|C_0|$ and the same vector-space dimension $d+1$ \textup(gcd-trivial because a new common root in $H\setminus C_0$ would have been a common root of $W$\textup).  The first part applies whenever $d\le j-|C_0|$ and gives $\binom{|H|-|C_0|}d$.
\end{proof}

\begin{remark}[what remains open]\label{rem:v13-conjf-open}
\Cref{thm:v13-fixeddim} shows that the primitive Conjecture-F core is a dimension-growing problem: common-root, quotient-periodic, dimension-one, and projective-plane strata are paid by \cref{lem:v13-gcd,lem:v13-quot-pullback,lem:v13-dim1,thm:v13-dim2}, and every fixed-dimensional stratum is paid by \cref{thm:v13-fixeddim}, but the growing-dimensional incidence theorem needed for the aperiodic band remains open and is part of \cref{prob:band}.
\end{remark}

\subsection{Quotient seam arithmetic and Hankel tools}

\begin{lemma}[GAP-2 quotient seam]\label{lem:v13-gap2}
Let $n,k,A$ be fixed and put $j=n-A$, $t=A-k$, and $r=n-k$, so $j+t=r$.  If $M\mid\gcd(n,k)$, then $M\mid r$, and hence
\[
        M\mid j \quad\Longleftrightarrow\quad M\mid t.
\]
On such buckets all quotient parameters $n/M,k/M,A/M,j/M,t/M$ are integral.  If $M\mid n$ and $M\mid j$ but $M\nmid k$, then the support descends but the dimension and window do not: a non-rate-preserving seam.
\end{lemma}

\begin{proof}
Since $M\mid n$ and $M\mid k$, we have $M\mid n-k=r=j+t$, so $j\equiv-t\pmod M$, proving the equivalence, and then $A=n-j$ and $t$ descend whenever $j$ does.  If $M\mid n$ and $M\mid j$ but $M\nmid k$, then $A=n-j$ is divisible by $M$ while $k/M$ and $t/M=(A-k)/M$ are not integral.
\end{proof}

For a finite ordered set $X$ of nonzero field elements and $m\ge1$, write $V_m(X)$ for the matrix with rows indexed by $x\in X$, columns $0\le r<m$, and entries $x^r$.

\begin{proposition}[Hankel factorization]\label{prop:v13-hankel}
Let $H\subseteq F^\times$ be finite.  For a word $w:H\to F$, define the rectangular syndrome Hankel block
\[
        H_{t,s}(w)_{a,b}:=\sum_{x\in H}w(x)x^{a+b}
        \qquad(0\le a<t,\ 0\le b<s).
\]
Then
\[
        H_{t,s}(w)=V_t(H)^{\mathsf T}\operatorname{diag}\bigl(w(x)\bigr)V_s(H).
\]
If $X=\{x_0,\ldots,x_{m-1}\}$ is a set of distinct nonzero points and
\[
        A_h(X):=V_m(X)^{\mathsf T}\operatorname{diag}(x^h:x\in X)\,V_m(X),
\]
then
\[
        \det A_h(X)=\det\bigl(V_m(X)\bigr)^2\prod_{x\in X}x^h\ \ne0 .
\]
\end{proposition}

\begin{proof}
The $(a,b)$ entry of the triple product is $\sum_{x}x^aw(x)x^b$, the defining Hankel entry.  The determinant identity is multiplicativity of the determinant together with $\det(V_m(X)^{\mathsf T})=\det V_m(X)$, and the Vandermonde determinant on distinct points is nonzero.
\end{proof}

\begin{proposition}[Johnson exchange mixing]\label{prop:v13-johnson-exchange}
Let $1\le j\le n-1$, and let $J(n,j)$ be the Johnson graph on $j$-subsets of an $n$-set, with degree $d=j(n-j)$.  If $\mathcal A$ is a family of $j$-subsets of density $\delta$ and $T_1$ is obtained from a uniform $T_0$ by one uniform exchange step, then
\[
        \Pr[T_0\in\mathcal A,\ T_1\in\mathcal A]
        \le
        \delta^2+\Bigl(1-\frac{n}{d}\Bigr)\delta(1-\delta),
\]
and point dictators $\mathcal A=\{T:x_0\in T\}$ attain equality.
\end{proposition}

\begin{proof}
Let $P$ be the normalized adjacency operator of $J(n,j)$ and write $1_{\mathcal A}=\delta\mathbf 1+f$ with $f\perp\mathbf1$ and $\|f\|_2^2=\delta(1-\delta)$ under the uniform measure.  The eigenvalues of the Johnson graph are classically
\[
        \theta_i=(j-i)(n-j-i)-i,
        \qquad 0\le i\le \min(j,n-j),
\]
the association-scheme spectrum of $J(n,j)$.  For $i<\min(j,n-j)$,
\[
        \theta_i-\theta_{i+1}
        =(j-i)(n-j-i)-(j-i-1)(n-j-i-1)+1
        =n-2i>0,
\]
so the sequence is strictly decreasing and the largest nontrivial normalized eigenvalue is $\lambda_1=\theta_1/d=1-n/d$.  Hence
\[
        \Pr[T_0,T_1\in\mathcal A]
        =\langle1_{\mathcal A},P1_{\mathcal A}\rangle
        =\delta^2+\langle f,Pf\rangle
        \le\delta^2+\lambda_1\|f\|_2^2 ,
\]
which is the display; only the largest nontrivial eigenvalue is needed for this upper bound.  The centered indicator of a point dictator lies in the $i=1$ eigenspace, so equality is attained there.
\end{proof}

\begin{proposition}[anticode packing]\label{prop:v13-anticode}
Let $0\le s\le j\le n$.  If a family $\mathcal F$ of $j$-subsets of an $n$-set satisfies $|A\setminus B|>s$ for all distinct $A,B\in\mathcal F$, then
\[
        |\mathcal F|\binom js\le \binom n{j-s}.
\]
\end{proposition}

\begin{proof}
For each $A\in\mathcal F$, consider its $(j-s)$-subsets, $\binom j{j-s}=\binom js$ of them.  These collections are disjoint across $\mathcal F$: a common $(j-s)$-subset of $A\ne B$ would give $|A\cap B|\ge j-s$, hence $|A\setminus B|\le s$, contrary to hypothesis.  There are at most $\binom n{j-s}$ subsets of size $j-s$ in total.
\end{proof}

\subsection{SPI deficiency-one chart}

Let
\[
        M(Z)=M_0+ZM_1
\]
be a $t\times(j+1)$ matrix pencil with affine-linear entries over $F$, and assume $t=j$.  For $0\le i\le j$, let $M_i(Z)$ be the maximal minor obtained by deleting column $i$, put
\[
        c_i(Z)=(-1)^iM_i(Z),
        \qquad
        L_Z(X)=\sum_{i=0}^{j}c_i(Z)X^i,
\]
and recall that whenever $\operatorname{rank}M(z)=j$, the cofactor vector $(c_0(z),\ldots,c_j(z))$ is nonzero and spans $\ker M(z)$.  In the syndrome application, $M(Z)$ is the exact-agreement syndrome pencil of \cref{sec:aperiodic-hankel-certificates} in the deficiency-one shape $t=j$, the kernel vector is the coefficient vector of a candidate locator of degree $\le j$, and a slope $z$ \emph{passes the split test} if the normalized kernel locator divides $X^n-1$; as in \cref{lem:singular-pencil}, passing the split test is necessary for badness at the exact agreement but not sufficient.

\begin{theorem}[deficiency-one eliminant or residual]\label{thm:v13-spi}
Assume $H$ is the full root set of $X^n-1$ and $\operatorname{char}F\nmid n$.
\begin{enumerate}[label=\textup{(\arabic*)},start=0]
\item If every maximal minor of $M(Z)$ vanishes identically, the pencil is identically rank-deficient; this chart is a named residual branch, to be classified by the polynomial-kernel chain of \cref{lem:singular-pencil}, and the statements below do not apply to it.
\end{enumerate}
Otherwise fix a maximal minor that is not identically zero.  Every finite slope $z$ then lies in exactly one of the classes:
\begin{enumerate}[label=\textup{(\roman*)}]
\item rank drop: all maximal minors vanish at $z$; this set is contained in the root set of the chosen minor, a nonzero polynomial of degree at most $t$;
\item low-degree full-rank chart: $\operatorname{rank}M(z)=j$ and $c_j(z)=0$; the Cramer locator has degree $<j$ and is charged to a contained higher-agreement branch;
\item top chart: $c_j(z)\ne0$; the Cramer locator is proportional to a monic degree-$j$ locator, and $z$ passes the split test if and only if $L_z(X)\mid X^n-1$.
\end{enumerate}
If $c_j\equiv0$, the top chart is empty and every full-rank slope lies in class \textup{(ii)}.  Assume $c_j\not\equiv0$ and perform pseudo-division over $F[Z]$:
\[
        c_j(Z)^\Delta\,(X^n-1)=Q(Z,X)\,L_Z(X)+R(Z,X),
        \qquad
        \Delta=n-j+1,
\]
with $\deg_XR<j$ and $\deg_ZR_m\le(n-j+1)t$ for every $X$-coefficient $R_m(Z)$ of $R$.  Then exactly one of the following holds:
\begin{enumerate}[label=\textup{(\alph*)}]
\item some $R_m$ is not identically zero; then the polynomial $c_j(Z)R_m(Z)$, of degree at most
\[
        t+(n-j+1)t,
\]
is a finite eliminant whose root set contains every top-chart slope passing the split test together with every slope of classes \textup{(i)} and \textup{(ii)};
\item all $R_m$ vanish identically; then every top-chart slope passes the split test.  Moreover, after monic normalization the Cramer locator is a fixed split divisor of $X^n-1$.  In the exact syndrome-pencil application this full-rank top chart is a contained fixed-support chart and contributes no support-wise noncontained exact-$A$ slope through this chart; roots of $c_j$ remain in the rank-drop ledger.
\end{enumerate}
\end{theorem}

\begin{proof}
Class \textup{(0)} and the trichotomy are immediate from the definitions: each maximal minor is the determinant of a $j\times j$ matrix with affine-linear entries in $Z$, hence a polynomial of degree at most $j=t$; if all vanish at $z$ then $\operatorname{rank}M(z)<j$, and otherwise the kernel is one-dimensional and spanned by the nonzero cofactor vector, whose top coordinate $c_j(z)$ decides between \textup{(ii)} and \textup{(iii)}.  In class \textup{(iii)}, $L_z$ has degree exactly $j$; since $X^n-1$ is squarefree with root set $H$, its monic normalization is $H$-split if and only if it divides $X^n-1$, which is the split test.

For the pseudo-division, dividing the degree-$n$ polynomial $X^n-1$ by the degree-$j$ polynomial $L_Z$ with leading coefficient $c_j(Z)$ performs $\Delta=n-j+1$ elimination steps, each multiplying the running remainder by $c_j(Z)$ once and subtracting a $c_i(Z)$-multiple of a shift of $L_Z$; since every $c_i$ has $Z$-degree at most $t$, each step raises the $Z$-degree of the remainder coefficients by at most $t$, giving $\deg_ZR_m\le\Delta t=(n-j+1)t$ from the degree-$0$ start.

In case \textup{(a)}, let $z$ be a top-chart slope passing the split test.  Substituting $Z=z$, the left-hand side $c_j(z)^\Delta(X^n-1)$ is divisible by $L_z$, and so is $Q(z,X)L_z(X)$; hence $L_z(X)\mid R(z,X)$, and $\deg_XR(z,X)<j=\deg L_z$ forces $R(z,X)=0$, so $R_m(z)=0$ for every $m$, in particular for the chosen nonzero one.  Slopes of classes \textup{(i)} and \textup{(ii)} satisfy $c_j(z)=0$ \textup(at rank drop, all minors vanish, in particular $M_j$\textup).  Hence all these slopes are roots of $c_jR_m$, whose degree is at most $t+(n-j+1)t$.

In case \textup{(b)}, work over $K=F(Z)$.  Since $c_j\not\equiv0$, the monic normalization $\widehat L_Z=L_Z/c_j$ is a degree-$j$ divisor of $X^n-1$ in $K[X]$.  The polynomial $X^n-1=\prod_{h\in H}(X-h)$ is split and squarefree over the constant field $F$, so every monic degree-$j$ divisor over $K$ is the product of a fixed $j$-subset of these linear factors.  Hence $\widehat L_Z=L_T$ for some fixed $T\subset H$, $|T|=j$, and $L_Z=c_jL_T$ in $F[Z,X]$.  The cofactor identity gives
\[
        M(Z)(c_0(Z),\ldots,c_j(Z))^{\mathsf T}=0,
\]
so, writing $\ell_T$ for the coefficient vector of the monic locator $L_T$,
\[
        c_j(Z)M(Z)\ell_T=0.
\]
Because $F[Z]$ is a domain and $c_j\not\equiv0$, one has $M(Z)\ell_T=0$.  In the exact syndrome-pencil application $M(Z)=U+ZV$, with $U=H_{j,j}(\operatorname{Syn}(f))$ and $V=H_{j,j}(\operatorname{Syn}(g))$, so $U\ell_T=V\ell_T=0$.  By the support-locator recurrence and the exact line-image map, both line coordinates are explained on $H\setminus T$, and the support-wise noncontainment condition $V\ell_T\ne0$ fails.  Thus the full-rank top chart is contained.  If $c_j(z)=0$ in this branch, then $c_i(z)=c_j(z)\ell_{T,i}=0$ for all $i$, so such slopes lie in the rank-drop ledger.
\end{proof}

\begin{remark}[relation to the v12 eliminant machinery, and calibration]\label{rem:v13-spi-calibration}
\Cref{thm:v13-spi} complements the certificate framework of \cref{sec:aperiodic-hankel-certificates}: \cref{thm:aperiodic-affine-eliminant} and \cref{thm:kernel-compressed-affine-eliminant} bound chart contributions \emph{given} a nonzero eliminant, while \cref{thm:v13-spi} \emph{produces} one for every deficiency-one chart with explicit degree, or else names the identically rank-deficient residual branch handled by \cref{lem:singular-pencil}.  The identically split full-rank top chart is not a support-wise noncontained residual after the exact line-image gate is imposed.  As calibration: for $n=512$, $k=256$, $A=384$, one has $t=j=128$, and the global one-parameter cap of case \textup{(a)} is
\[
        t+(n-j+1)t=128+385\cdot128=49408,
\]
a bounded eliminant problem.  The theorem does not classify higher-deficiency SPI charts; those remain inside \cref{prob:band}.
\end{remark}

\section{v13 row-packet schema and non-claims}\label{sec:v13-packet-schema}

A row packet invoking this insert should print enough data for an independent integer check:
\[
\begin{array}{ll}
\text{row} & (F,D,k,n,\rho),\\
\text{denominators} & q_{\mathrm{gen}},\,q_{\mathrm{line}},\,q_{\mathrm{chal}},\,q_{\mathrm{list}},\\
\text{target} & \varepsilon^*,\quad B_*=\floor{\varepsilon^*Q},\\
\text{agreement interval} & I=[a_{\min},a_{\max}]\cap\mathbb Z,\\
\text{unsafe certificates} & L(a)>B_*,\\
\text{safe certificates} & U(a)\le B_*,\\
\text{paid cells} & \Paid_{\tan},\ \Paid_{\mathrm{quot}},\ \Paid_{\mathrm{ext}},\\
\text{residual cells} & \text{named aperiodic, extension, list, or SPI branch},\\
\text{deduplication rule} & \text{support/image/root coalescing theorem used},\\
\text{endpoint convention} & \text{closed integer ball and real supremum \textup(\cref{cor:v13-endpoint}\textup)}.
\end{array}
\]
This insert does not prove an aperiodic M1 local limit, a primitive L1 image-fiber upper theorem, a new extension-line transfer theorem, a protocol soundness theorem, or a new field-size cap.  Its role is to make those future inputs auditable once supplied: the tangent cell is exact only in the printed deep Reed--Solomon range $3(n-A)\le n-k$; quotient and extension cells must be paid with explicit certificates at the printed status grades; and residual branches must remain named rather than hidden inside point estimates.

\section{v13 integration checklist}\label{sec:v13-checklist}

A CAP25 v13 update can use this insert as follows.

\begin{enumerate}[label=\textup{(\arabic*)},leftmargin=2.2em]
\item In the threshold section, cite \cref{thm:v13-staircase,cor:v13-endpoint,thm:v13-corridor,thm:v13-steepness,thm:v13-windows} for the exact integer certificate grammar; \cref{prop:onestep} is subsumed as the one-step case \textup(\cref{rem:v13-staircase-v12}\textup).  Use \cref{def:v13-quotient-status} to prevent collision-open quotient zones from being printed as point estimates.
\item Use \cref{prop:v13-tangent,prop:v13-quotient-safe-sum,prop:v13-extension} as paid-cell compilers.  The tangent term is exact in the range $3(n-A)\le n-k$ by \cref{prop:v13-tangent} with upper bound \cref{thm:deep-mca}; the quotient cell cites the v12 ledgers at its status grade; the extension-pole term requires the stated conversion and pairwise pole-collision certificate, realized in v12 by \cref{cor:extension-pole-deep-list-floor}.
\item Use \cref{thm:v13-planted,cor:v13-polyfold-planted,cor:v13-list-unsafe,prop:v13-dyadic-planted} for lower-side list caps from planted quotient-core or polynomial-folding rows, and \cref{cor:v13-list-windows} for exact unresolved denominator windows.
\item Use \cref{thm:v13-johnson-list} to close ordinary Johnson-region list tails; for interleaved rows the arity-uniform \cref{thm:johnson-list} applies unchanged.
\item Use \cref{prop:v13-concrete-sunflower,prop:v13-pma,cor:v13-pma-johnson,prop:v13-zero-excess-chart,prop:v13-distinct-excess-chart,prop:v13-fixed-excess} as paid compiler pieces inside the L1 sunflower program, and state \cref{prob:v13-primitive-image-fiber,prob:v13-l1-residuals} as the remaining L1 obstructions.
\item Use \cref{lem:v13-gcd,lem:v13-quot-pullback,lem:v13-gap2} to remove paid common-divisor and quotient-periodic strata before calling a branch aperiodic in the sense of \cref{def:periodicity-scale}.
\item Use \cref{thm:v13-fixeddim,thm:v13-dim2,lem:v13-dim1} to dispose of fixed-dimensional and projective-plane Conjecture-F subcases.  The growing-dimensional primitive incidence problem remains open inside \cref{prob:band}.
\item Use \cref{thm:v13-first-moment,thm:v13-second-moment,cor:v13-exact-variance,cor:v13-dependency} as a random-pair moment ledger only, never as a worst-case safe-side theorem \textup(\cref{rem:v13-moment-scope}\textup).
\item Use \cref{prop:v13-hankel,prop:v13-johnson-exchange,prop:v13-anticode} as algebraic and support-combinatorial inputs for exchange-rigidity or spectral routes.
\item Use \cref{thm:v13-spi} to state the first SPI chart as eliminant-or-contained-chart: nonzero pseudo-remainders give finite eliminants, identically rank-deficient pencils remain a named residual branch connected to \cref{lem:singular-pencil}, and identically split full-rank top charts are contained after the exact line-image noncontainment gate.
\end{enumerate}

% ============================= END V13 BODY =============================
